\documentclass[12pt]{article}

\usepackage[a4paper, margin=1in]{geometry}
\usepackage{amsmath,amssymb,amsthm}
\usepackage{mathtools}
\usepackage{mathrsfs}
\usepackage{enumitem}
\usepackage{booktabs}
\usepackage{array}
\usepackage{xcolor}
\usepackage[hidelinks]{hyperref}
\usepackage{fontspec}

\defaultfontfeatures{Ligatures=TeX}
\IfFontExistsTF{Times New Roman}
  {\setmainfont{Times New Roman}}
  {\setmainfont{TeX Gyre Termes}}
\IfFontExistsTF{Menlo}
  {\setmonofont{Menlo}}
  {\setmonofont{Latin Modern Mono}}

\newtheorem{theorem}{Theorem}[section]
\newtheorem{proposition}[theorem]{Proposition}
\newtheorem{lemma}[theorem]{Lemma}
\newtheorem{corollary}[theorem]{Corollary}
\newtheorem{conjecture}[theorem]{Conjecture}
\theoremstyle{definition}
\newtheorem{definition}[theorem]{Definition}
\newtheorem{example}[theorem]{Example}
\theoremstyle{remark}
\newtheorem{remark}[theorem]{Remark}

\newcommand{\Fix}{\operatorname{Fix}}
\newcommand{\Hull}{\operatorname{Hull}}
\newcommand{\Stab}{\operatorname{Stab}}
\newcommand{\Vir}{\mathrm{Vir}}
\newcommand{\AU}{\mathcal A_U}
\newcommand{\Ad}{\operatorname{Ad}}

\title{%
  \textbf{Foundation I: From Source to Manifestation}\\[0.4ex]
  \large\textbf{Accessible Algebras from a Source,}\\
  \large\textbf{Its Noise and a Readable Seed}%
}

\author{
  Sidong Liu, PhD \\
  \small iBioStratix Ltd \\
  \small \texttt{sidongliu@hotmail.com}
}

\date{October 2026}
\hypersetup{
  pdftitle={Foundation I: From Source to Manifestation. Accessible Algebras from a Source, Its Noise and a Readable Seed},
  pdfauthor={Sidong Liu}
}

\begin{document}
\emergencystretch=2em
\raggedbottom
\hbadness=5000
\vbadness=5000

\maketitle

\begin{abstract}
A source, its noise and a physical reader supply different constraints on the
structure in which a world can be read. We separate finite measurement permissions
from the ideal algebra generated by their effects. A partially mixed record probe
generates the same pointer algebra at every nonzero gain, although its sharp-test
error approaches one half as the gain vanishes. Repeated fresh nondemolition probes
give a finite instrument-level approximation to the sharp measurement, including
old references; this supplies a resource-sensitive bridge to finite tasks.

For a stationary source with a faithful state, state-preserving noise has a
modular-invariant fixed algebra and a unique state-preserving expectation. Its
Ces\`aro projection need not be a long-time limit of individual trajectories.
A stable seed has a modular hull inside that fixed algebra. If the ambient source
has positive translations with only the vacuum invariant, every nonscalar hull
closed under forward translations inherits a proper Cut on its vacuum space.
A minimal compact symmetry identifies any hull containing the symmetry-fixed
algebra with the fixed algebra of its seed stabilizer, provided the closure maps
preserve it. In the SU(2) level-one net, a positive log-bump stress seed
generates the Virasoro half-line algebra under modular evolution alone and
its one-sided orbit forms a standard Cut on the Virasoro vacuum space.
One bounded current seed generates the
current algebra under modular flow alone; its character distinguishes the
stabilizer up to conjugacy. Axis dephasing leaves three hulls for the stated
stress/current seeds, and weak symmetry breaking determines their passive
lifetimes. Stress excitations have the same relative entropy across these hulls,
whereas current excitations distinguish them. These results derive the source--access--Cut chain in its net-theoretic geometry. The ambient net, available couplings and preparation permissions remain
explicit inputs; forward translations are unnecessary for the stated
log-bump stress seeds.
\end{abstract}

\clearpage
\tableofcontents
\clearpage

\section{Introduction}\label{sec:intro}

Consciousness~I~\cite{PaperI} works with an accessible algebra \(A\) inside a global
algebra \(M\) with vacuum \(\Omega\). When \(A\) contains a Cut, a subalgebra stable
under the modular evolution in one direction only, the translations of a chiral
conformal net are computed from modular data, and for a standard Cut the net and the
modular data determine each other. In the SU(2) level-one current net the same paper
finds that the algebra fixed by an explicit rotation noise, the \(c=1\) Virasoro
algebra of a half-line, coincides with the algebra generated from one bounded stress
observable by dilations and forward translations. For a positive log-coordinate
bump with monotone ends, the construction needs only modular evolution:
its one-sided orbit forms a proper standard Cut on the Virasoro vacuum space,
and its full modular hull is the Virasoro half-line algebra. The origin of the
ambient structure and of physically available seed operations remains open;
the present paper asks which accessible algebra a source and seed select.

Foundation~I is the source end of that chain. Its question is what an actual source,
physical access to it and persistent organization fix on the way to an accessible
algebra, and what they leave free. The HAFF papers supply part of the answer in finite
dimension. Paper~A~\cite{HAFFA} shows that two channels can generate the same algebra
and still resolve different things, so that one state and one reconstruction rule give
different geometries. Paper~B~\cite{HAFFB} derives the exact fixed algebra in the commuting-coupling
class and the unique largest block algebra passing its stated window criteria.
Block-constant recording couplings expose central labels, while protected
relations, classical sector correlations and record-conditioned formation have
different roles. Its collective recording and Hamiltonian-link constructions
show how actual process resources create relations; the noise graph alone does
not determine all correlations. Paper~E~\cite{HAFFE}
separates what a probe can implement, what the dynamics leaves fixed and what fragments
record, and computes all three in one collision model. For its Gaussian
linear-response task with quadratic loss and an unrestricted rank-\(k\)
projector menu, T-DOME~II~\cite{TDOMEII} identifies a top-\(k\)
posterior-task eigenspace as an optimal reading subspace. A restricted
physical instrument menu need not realize that optimizer. What these results do not say is how a seed obtained from access
grows into an algebra that can carry geometry, when the result is proper, how it
depends on the seed, and whether different seeds give different geometries. HAFF-D
lists this as the missing selection step: with a faithful ambient state the ambient
algebra itself passes every modular and semigroup stability test, and a nontrivial
selection needs a seed and a minimal hull~\cite{HAFFD}.

The objects are as follows. A \emph{source} is a von Neumann algebra \(\AU\) with a
cyclic and separating vector \(\Omega\), its state \(\omega\), and a semigroup
\(\{\mathcal E_t\}\) of normal unital completely positive maps describing the
environment in the Heisenberg picture. The modular group \(\sigma_s\) of \(\omega\) is
fixed by the source and needs no further input; for a thermal state it is the time
evolution in units of the inverse temperature, and for the vacuum of a wedge it is the
boost group. \emph{Access} supplies actual effects through physical probes; a reader may keep a
set of them as a \emph{seed}. Algebraic functional calculus then describes the
ideal distinctions encoded by that seed. It is a further closure, not an
additional measurement permission. The
\emph{stable hull} of the seed is the smallest von Neumann algebra that contains it and
is closed under the modular flow and environmental dynamics. When an ambient positive translation structure is
available, we also study its closure under forward translations. These are distinct
closures: translation stability must be imposed or proved before a Cut can be inherited.

The argument proceeds from physical access to persistent structure:
\begin{enumerate}
\item Section~\ref{sec:permission} gives the finite pointer-instrument
  bridge, including old references and adaptive task error.
\item Sections~\ref{sec:stationary}--\ref{sec:hull} derive the
  state-preserving expectation and modular stable hull from
  stationarity, and separate this projection from dynamical relaxation.
\item Section~\ref{sec:cut} proves proper Cut inheritance for
  translation-stable nonscalar hulls with unique invariant vacuum.
\item Section~\ref{sec:symmetry} identifies the classified hull by its
  seed stabilizer. Current and positive log-bump stress seeds admit
  modular-only generation; characters distinguish stabilizers and
  symmetry-breaking noise fixes passive lifetimes.
\item Sections~\ref{sec:foreground}--\ref{sec:response} separate
  ideal seed closure from finite coding budgets and distinguish
  seed-independent stress response from seed-dependent current response.
\end{enumerate}

The ambient algebra, its translations, the coupling to the environment and the task
that fixes the foreground remain inputs; Section~\ref{sec:open} lists them with the
conjectures that would remove some of them. The finite-dimensional statements are
checked by scripts described in Appendix~\ref{app:checks}.

\section{Source, access, seed and hull}\label{sec:objects}

\begin{definition}[Source]\label{def:source}
A \emph{source} is a quadruple \((\mathcal H_U,\AU,\Omega,\{\mathcal E_t\}_{t\ge0})\):
a von Neumann algebra \(\AU\) on \(\mathcal H_U\), a vector \(\Omega\) cyclic and
separating for \(\AU\), and a semigroup of normal unital completely positive maps
\(\mathcal E_t:\AU\to\AU\), pointwise continuous in the \(\sigma\)-weak topology. We write
\(\omega=\langle\Omega,\cdot\,\Omega\rangle\), \(\Delta,J\) for the modular operator
and conjugation of \((\AU,\Omega)\), \(\sigma_s=\Ad\Delta^{is}\) for the modular group,
and
\[
\Fix(\mathcal E)=\{x\in\AU:\mathcal E_t(x)=x\ \text{for all}\ t\ge0\}.
\]
The source is \emph{stationary} if \(\omega\circ\mathcal E_t=\omega\) for all \(t\).
An \emph{ambient positive translation structure} (a Borchers pair with $\AU$) is a strongly continuous unitary group \(U(a)=e^{iaP}\) with
\(P\ge0\), \(U(a)\Omega=\Omega\) and \(U(a)\AU U(a)^*\subseteq\AU\) for \(a\ge0\).
\end{definition}

The ambient positive translation structure is the structure of a half-line algebra of a chiral net in its vacuum,
or of a wedge algebra with a lightlike translation. By Borchers' theorem it satisfies
\(\Delta^{it}U(a)\Delta^{-it}=U(e^{-2\pi t}a)\) and \(JU(a)J=U(-a)\)~\cite{Borchers1992},
and \(U(1)\AU U(1)^*\subset\AU\) is then a half-sided modular inclusion whose
translation group, in the sense of Wiesbrock's theorem, is \(U\)~\cite{Wiesbrock1993,AZ}. We use it only where it is stated.

\paragraph{Access and seed.} Paper~E of HAFF describes access by a context: allowed
probe channels \(\mathcal R_\alpha^*:\mathcal B(\mathcal K_\alpha)\to\AU\), whose images
of probe effects are the implementable effects \(\mathrm{Eff}\), the environmental
dynamics with its fixed set, and a record algebra~\cite{HAFFE}. Paper~B identifies central record labels for its block-constant recording
couplings, not for every possible environment instrument~\cite{HAFFB}.
A seed \(S\subseteq\AU\) is a set of bounded self-adjoint observables.
A \emph{physical effect seed} is selected from the actual probe effects and
retained under a declared budget (Section~\ref{sec:foreground}); an
\emph{ideal spectral seed} may additionally contain bounded Borel functions of
those effects. Both enter the same mathematical hull, but the latter need not
be implementable by the former probe menu. A seed is \emph{stable} if
\(S\subseteq\Fix(\mathcal E)\). The hull definitions do not depend on the
mechanism selecting the seed. Section~\ref{sec:permission} supplies a nonempty
physical bridge and exhibits the information lost by identifying permissions
with this ideal closure.

\begin{definition}[Stable hull]\label{def:hull}
Let \(\Gamma\) be a family of normal maps of \(\AU\) into itself that contains the
modular group, \(\{\sigma_s\}_{s\in\mathbb R}\subseteq\Gamma\). A von Neumann
subalgebra \(\mathcal N\subseteq\AU\) is \emph{\(\Gamma\)-admissible} if
\(\gamma(\mathcal N)\subseteq\mathcal N\) for every \(\gamma\in\Gamma\). The
\emph{\(\Gamma\)-hull} of a seed \(S\) is
\[
\Hull_\Gamma(S)=\bigcap\{\mathcal N:\ \mathcal N\ \Gamma\text{-admissible},\ S\subseteq\mathcal N\}.
\]
We use three families: \(\Gamma_\sigma=\{\sigma_s\}\) (the \emph{modular hull}),
\(\Gamma_2=\{\sigma_s\}\cup\{\mathcal E_t\}\) (the \emph{stable hull}), and, in the
presence of an ambient positive translation structure,
\(\Gamma_3=\Gamma_2\cup\{\Ad U(a)\}_{a\ge0}\)
(the \emph{translation-stable hull}). Existence of an ambient positive translation structure alone does not
identify the \(\Gamma_2\)- and \(\Gamma_3\)-hulls.
\end{definition}

Since \(\sigma_{-s}\in\Gamma\), an admissible algebra satisfies
\(\sigma_s(\mathcal N)=\mathcal N\). For a seed of exact pointer observables
realising a foreground frame, \(\Hull_{\Gamma_2}\) is the minimal modular--decoherence
hull of that frame. The \((\sigma,U)\)-hull of Consciousness~I is \(\Hull_{\Gamma_3}\) without the noise
maps~\cite{PaperI}. HAFF-D calls the \(\Gamma_2\)-admissible algebras modular-stable
geometric candidates and notes that the class contains the ambient algebra, so that a
selection needs a seed~\cite{HAFFD}.

\section{From a physical effect to a finite task}\label{sec:permission}

A von Neumann algebra closes under bounded spectral functions, products and
operator limits. Probe permissions need not close under these operations.
Even the linear inverse of a noisy effect may require a negative coefficient,
so it cannot in general be performed by randomizing permitted tests.

\begin{proposition}[Same algebra, different sharp permissions]\label{prop:unsharp}
Let \(P_0,P_1\) be the pointer projections of a qubit. A fresh memory
\(\eta_\alpha=(\mathbf1+\alpha Z)/2\), \(0\le\alpha\le1\), interacts through
\[
V=P_0\otimes\mathbf1+P_1\otimes X .
\]
A \(Z\) readout of the memory implements
\[
F_\alpha=\tfrac{1+\alpha}{2}P_0+\tfrac{1-\alpha}{2}P_1.
\]
For \(\alpha>0\), \(W^*(F_\alpha)=W^*(P_0)\), and every hull in
Definition~\ref{def:hull} gives the same result for these two seeds.
With arbitrary memory effects but only this single collision, the closest
implementable effect to \(P_0\) has operator-norm error
\((1-\alpha)/2\). Thus equality of the ideal hulls does not give equality
of their sharp-measurement permissions.
\end{proposition}

\begin{proof}
For \(\alpha>0\),
\(P_0=[F_\alpha-(1-\alpha)\mathbf1/2]/\alpha\); each seed belongs to the
algebra of the other. Paper~E computes the exact single-port permission set
as the parallelogram with vertices
\((0,0),(1,1),((1+\alpha)/2,(1-\alpha)/2)\) and its transpose,
where \((a,b)\) denotes \(aP_0+bP_1\)~\cite{HAFFE}.
In that set \(a-b\le\alpha\), so
\(\max(1-a,b)\ge(1-\alpha)/2\); \(F_\alpha\) attains the bound.
At \(\alpha=0\) all effects are scalar and the seed algebra is trivial.
\end{proof}

All nonzero gains therefore have the same ideal structure, while their
finite-read abilities approach the uninformative probe continuously as
\(\alpha\downarrow0\). The spectral inverse is an algebraic identity.
It neither sharpens a one-shot probe nor specifies its back-action.

There is an abstract bounded coupling for an unsharp version of any
self-adjoint contraction \(O\). Put
\[
\vartheta(O)=\tfrac12\arccos(\alpha O),\qquad
V_O=\exp[-i\vartheta(O)\otimes Y].
\]
A pure qubit port prepared in \(|0\rangle\) and read in \(Z\) yields
\((\mathbf1+\alpha O)/2\). For a chosen gate duration \(t_g\),
\(H_O=(\hbar/t_g)\vartheta(O)\otimes Y\) is bounded, and adding
\(\pi\hbar\mathbf1/(2t_g)\) makes it positive without changing the instrument.
For \(\alpha>0\), this effect generates \(W^*(O)\); if \(O\) is fixed by
stationary noise, its spectral functions are fixed by
Theorem~\ref{thm:stationary}(a). This construction specifies the required
interaction, preparation and readout. It does not derive a local field
Hamiltonian implementing \(\vartheta(O)\), or grant that coupling to a
reader whose original menu lacks it. In particular, the current and stress
seeds below require a physical bounded-function readout or another justified
representation of that readout.

\begin{proposition}[Finite sharp-instrument approximation]\label{prop:readbridge}
Fix \(0<\alpha\le1\). Let \(m\ge1\) fresh memories of polarization \(\alpha\)
undergo the preceding controlled-X collision with the \emph{same} source
qubit. Read each memory in \(Z\), retain the majority label, and use a fair
independent tie breaker if \(m\) is even. No rotation of the source occurs
between collisions. Let \(e_m\) be the majority error of \(m\) independent
binary tests of error \((1-\alpha)/2\). After the fine records are discarded,
the coarse instrument, including its output source, is
\[
\mathcal I_m(\rho)=
\sum_{y=0,1}|y\rangle\langle y|\otimes
 \sum_{z=0,1}q_m(y|z)P_z\rho P_z,\qquad
q_m(z|z)=1-e_m,\quad q_m(1-z|z)=e_m .
\]
For the sharp L\"uders instrument
\(\mathcal I_{\mathrm{sh}}(\rho)=\sum_z|z\rangle\langle z|\otimes P_z\rho P_z\),
\[
\tfrac12\|\mathcal I_m-\mathcal I_{\mathrm{sh}}\|_\diamond=e_m
\le e^{-m\alpha^2/2}.
\]
This bound holds on arbitrary old references of the source.
\end{proposition}

\begin{proof}
The memory state is diagonal. After one collision, projection onto either
memory \(Z\) outcome removes the source off-diagonal blocks: the partial
maps are
\(\rho\mapsto\sum_z q_1(y|z)P_z\rho P_z\).
Fresh memories give the product probabilities conditional on the same
unchanged pointer label, as a consequence of this joint instrument.
Grouping all fine outcomes by their majority gives the displayed map.
For a joint input with a reference, put
\(R_z=(P_z\otimes\mathbf1)\rho(P_z\otimes\mathbf1)\).
The two output-label blocks of the difference are
\(e_m(R_1-R_0)\) and its negative. Their orthogonal source supports give
trace norm \(2e_m\operatorname{Tr}\rho\).
Equivalently the map is the sharp instrument followed by a binary label
flip, whose diamond distance from the identity is \(2e_m\); a pointer
input attains it. Hoeffding's inequality applied to the conditional
\(\{\pm1\}\) samples of mean \(\alpha\) gives the final bound.
\end{proof}

This is a physical finite route from an ideal pointer distinction to a
task permission. To reach error \(\epsilon>0\), it suffices to use
\(m\ge2\alpha^{-2}\log(1/\epsilon)\). The allowance must contain the
controlled-X coupling, fresh independently prepared memories, \(m\)
readouts, classical majority processing and storage; if orientation or
gain is unknown, Paper~E's paid references or another calibration are
also needed. Discarding the fine records means restricting the task
output channel; it does not physically reset their memories or erase
externally retained copies. The first read destroys pointer coherence exactly as does
the comparison sharp instrument. The claim concerns that instrument,
not undisturbed reading of an arbitrary unknown quantum state.

\begin{corollary}[Finite task transfer]\label{cor:tasktransfer}
Consider a finite adaptive circuit with arbitrary retained systems and
references. Suppose its \(q\) sharp pointer-instrument slots can each be
replaced by Proposition~\ref{prop:readbridge}, with half-diamond
errors at most \(e_i\), uniformly over earlier records. Keep all other
allowed operations the same in both circuits. Their output states have
trace distance at most \(\min(1,\sum_i e_i)\); a loss in \([0,L]\) changes
in expectation by at most \(L\min(1,\sum_i e_i)\).
\end{corollary}

\begin{proof}
Include the classical history in the retained state. Replacing one slot
at a time and using contractivity of all intervening physical channels
bounds the full diamond distance by \(2\sum_i e_i\). Trace distance
bounds the expectation of a loss in the stated interval.
\end{proof}

The same argument incorporates time-window or gate-error certificates
at their corresponding circuit slots; Paper~B gives such finite-time
risk bounds for its channel class~\cite{HAFFB}. Each slot requires an
instrument-level bound, not just nearby effects. A \(\Gamma\)-hull supplies
candidate ideal observables for a task. A physical compiler supplies
their actual channels and costs. The corollary connects these layers for
a finite protocol meeting its conditions; it does not make all modular
times, products, translations or von Neumann limits available at a
fixed seed budget.

For example, applying \(X\) to the output source after a sharp pointer
instrument leaves its effects unchanged but sends a \(P_0\) input to
\(P_1\). Its output is at trace distance one from the unmodified
instrument on that input. Equal probabilities, equal seed algebras and
even equal sharp effects do not certify the physical task after reading.

\section{Stationary sources}\label{sec:stationary}

Three structural conditions on the fixed set are usually imposed separately: that
\(\Fix(\mathcal E)\) be closed under products, that it be invariant under the modular
group, and that it carry an \(\omega\)-preserving conditional expectation. A
persistent Nakajima--Zwanzig projection is a fourth. For a stationary source all of
them follow.

\begin{theorem}[Stationary source]\label{thm:stationary}
Let the source be stationary. Then:
\begin{enumerate}[label=(\alph*),leftmargin=2em]
\item every \(a\in\Fix(\mathcal E)\) lies in the multiplicative domain of each
  \(\mathcal E_t\), so \(\mathcal E_t(axb)=a\,\mathcal E_t(x)\,b\) for
  \(a,b\in\Fix(\mathcal E)\), \(x\in\AU\), and \(\Fix(\mathcal E)\) is a von Neumann
  subalgebra;
\item there is a unique normal conditional expectation
  \(E_{\mathrm{fix}}:\AU\to\Fix(\mathcal E)\) with \(\omega\circ E_{\mathrm{fix}}=\omega\);
  it is faithful and
  \(E_{\mathrm{fix}}\circ\mathcal E_t=\mathcal E_t\circ E_{\mathrm{fix}}=E_{\mathrm{fix}}\);
\item \(\sigma_s(\Fix(\mathcal E))=\Fix(\mathcal E)\) for all \(s\), and the modular
  group of \(\omega|_{\Fix(\mathcal E)}\) is the restriction of \(\sigma\);
\item with \(\mathcal P=E_{\mathrm{fix}}\) and \(\mathcal Q=\mathrm{id}-\mathcal P\),
  \(\mathcal P\mathcal E_t\mathcal Q=\mathcal Q\mathcal E_t\mathcal P=0\) and
  \(\mathcal P\mathcal E_t\mathcal P=\mathcal P\): the projected dynamics on the fixed
  algebra is the identity, and the Nakajima--Zwanzig memory kernel between the fixed
  algebra and its complement vanishes.
\end{enumerate}
\end{theorem}

\begin{proof}
(a) For \(a\in\Fix(\mathcal E)\) the Kadison--Schwarz inequality gives
\(\mathcal E_t(a^*a)\ge\mathcal E_t(a)^*\mathcal E_t(a)=a^*a\), and
\(\omega(\mathcal E_t(a^*a)-a^*a)=0\) by stationarity. Since \(\omega\) is faithful,
\(\mathcal E_t(a^*a)=a^*a\); the same holds for \(aa^*\) because \(a^*\) is also fixed.
By Choi's theorem~\cite{Choi1974} \(a\) belongs to the multiplicative domain of
\(\mathcal E_t\), on which \(\mathcal E_t\) is a bimodule map. Products of fixed
elements are therefore fixed, and normality of \(\mathcal E_t\) makes the fixed set
\(\sigma\)-weakly closed.

(b) The idea is the mean ergodic theorem of Kov\'acs and Sz\H{u}cs~\cite{KovacsSzucs1966}; for semigroups of normal unital positive maps with a faithful invariant state see Watanabe~\cite{Watanabe1979}.
By Kadison--Schwarz and stationarity,
\(\|\mathcal E_t(x)\Omega\|^2\le\omega(\mathcal E_t(x^*x))=\|x\Omega\|^2\), so
\(V_tx\Omega=\mathcal E_t(x)\Omega\) defines a contraction semigroup on
\(\mathcal H_U\) with \(V_t\Omega=\Omega\). Its Ces\`aro means converge strongly to the
orthogonal projection \(P\) onto its fixed vectors. The means
\(x_T=T^{-1}\int_0^T\mathcal E_t(x)\,dt\) are bounded by \(\|x\|\) and satisfy
\(x_T\Omega\to Px\Omega\); every \(\sigma\)-weak cluster point \(y\) has \(y\Omega=Px\Omega\),
and since \(\Omega\) is separating the cluster point is unique, so \(x_T\to y\)
\(\sigma\)-weakly. Put \(E_{\mathrm{fix}}(x)=y\). Since
\(\|\mathcal E_s(x_T)-x_T\|\le2s\|x\|/T\), \(y\in\Fix(\mathcal E)\); on
\(\Fix(\mathcal E)\) the map is the identity, so \(E_{\mathrm{fix}}\) is a projection of
norm one onto a von Neumann algebra, hence a conditional expectation (Tomiyama). The expectation preserves the state because
\(\omega(E_{\mathrm{fix}}(x))=\langle\Omega,Px\Omega\rangle=\omega(x)\), as \(P\Omega=\Omega\).
To prove normality, let \(0\le x_i\uparrow x\) and
\(y=\sup_iE_{\mathrm{fix}}(x_i)\le E_{\mathrm{fix}}(x)\).
Normality of \(\omega\) gives
\[
\omega(y)=\sup_i\omega(E_{\mathrm{fix}}(x_i))
=\sup_i\omega(x_i)=\omega(x)=\omega(E_{\mathrm{fix}}(x)).
\]
Faithfulness forces \(y=E_{\mathrm{fix}}(x)\); thus the expectation is normal,
and state preservation makes it faithful. The means of
\(\mathcal E_t(x)\) differ from those of \(x\) by at most \(2t\|x\|/T\) in norm, which
gives \(E_{\mathrm{fix}}\circ\mathcal E_t=E_{\mathrm{fix}}\), and
\(\mathcal E_t\circ E_{\mathrm{fix}}=E_{\mathrm{fix}}\) because the range is fixed. A
normal \(\omega\)-preserving conditional expectation onto a given subalgebra is unique,
since it is determined by \(E(x)\Omega=e\,x\Omega\) with \(e\) the projection onto
\(\overline{\Fix(\mathcal E)\Omega}\)~\cite{Takesaki1972}.

(c) By Takesaki's theorem~\cite{Takesaki1972}, a von Neumann subalgebra carries a
faithful normal \(\omega\)-preserving conditional expectation if and only if it is
globally invariant under \(\sigma\), and then \(\sigma\) restricts to the modular group
of the restricted state.

(d) follows from (b).
\end{proof}

Stationarity supplies all four conditions. HAFF-E proves the
fixed-algebra part~\cite{HAFFE}; the vacuum-preserving internal noise
of Consciousness~I~\cite{PaperI} and the block-diagonal stationary
states of HAFF-B~\cite{HAFFB} give concrete instances.

What fails without stationarity is (c). If a faithful normal state \(\varphi\neq\omega\)
is invariant, (a) and (b) hold with \(\varphi\) in place of \(\omega\) and
\(\Fix(\mathcal E)\) is invariant under the modular group of \(\varphi\), which need not
preserve it for \(\omega\). An entangled faithful state of \(M_2\otimes M_3\) with noise
acting ergodically on the second factor is an example (check~3 of
Appendix~\ref{app:checks}): the fixed algebra is \(M_2\otimes\mathbf 1\), and the
modular group of the entangled state moves it.

\begin{example}[A non-stationary qubit]\label{ex:qubit}
Let \(\AU=M_2\), let \(\mathcal E_t\) multiply the off-diagonal entries in the
\(Z\)-basis by \(e^{-\gamma t}\), and let \(\omega=\operatorname{Tr}(\rho\,\cdot)\) with
\(\rho\) faithful and not diagonal. Then \(\Fix(\mathcal E)\) is the diagonal algebra,
but \(\sigma_s(P_0)=\rho^{is}P_0\rho^{-is}\) is a rank-one projection different from
\(P_0=|0\rangle\langle0|\) for small \(s\neq0\). Two distinct nonorthogonal rank-one projections
generate \(M_2\), so the hull of the single pointer projection \(P_0\) for any family
\(\Gamma\supseteq\Gamma_\sigma\) is all of \(M_2\). For the diagonal state
\(\rho_\infty=\operatorname{diag}(\rho)\), which is stationary, the hull of \(P_0\) is
the two-dimensional algebra it generates.
\end{example}

An invariant expectation supplies a stationary partner with the same values
on the persistent sector, whether or not individual trajectories relax.

\begin{proposition}[Stationary partner]\label{prop:asymptotic}
Suppose there is a faithful normal conditional expectation \(E_\infty\) of \(\AU\) onto
\(\Fix(\mathcal E)\) with \(E_\infty\circ\mathcal E_t=E_\infty\). Then
\(\omega_\infty=\omega\circ E_\infty\) is faithful, normal and stationary, agrees with
\(\omega\) on \(\Fix(\mathcal E)\), and Theorem~\ref{thm:stationary} applies to
\((\AU,\omega_\infty,\mathcal E)\).
\end{proposition}

\begin{proof}
\(\omega_\infty\circ\mathcal E_t=\omega\circ E_\infty\circ\mathcal E_t=\omega_\infty\);
faithfulness and normality are inherited from \(\omega\) and \(E_\infty\); on
\(\Fix(\mathcal E)\), \(E_\infty\) is the identity. In standard form \(\omega_\infty\) is
a vector state of a cyclic and separating vector.
\end{proof}

For the commuting couplings of HAFF-B, \(E_\infty\) is the pinching onto the blocks; for
the noises of Consciousness~I it is the Haar average. The partner state defines its own modular hull.
When these semigroups converge to their averages it is also the relaxed
state, but Proposition~\ref{prop:asymptotic} needs only the invariant
expectation. The change of hull in Example~\ref{ex:qubit} illustrates
why the partner and original source must remain distinguished.

\begin{example}[A mean without relaxation]\label{ex:mean}
On \(M_2\), let
\(\mathcal E_t(x)=e^{itZ}xe^{-itZ}\), with the faithful stationary
source state \(\tau=\mathbf1/2\). The fixed algebra is diagonal and its
Ces\`aro expectation is pinching. For
\(\rho=(\mathbf1+aX)/2\), \(0<a<1\), the Schr\"odinger trajectory is
\(e^{-itZ}\rho e^{itZ}\), whose mean tends to \(\mathbf1/2\), but
\[
\tfrac12\|e^{-itZ}\rho e^{itZ}-\mathbf1/2\|_1=a/2
\]
at every time. The fixed-sector projection in
Theorem~\ref{thm:stationary} is exact; it is not a claim that the
complementary dynamics decays.
\end{example}

\paragraph{Finite windows.} In its commuting class, HAFF-B selects a unique
largest admissible block algebra \(\mathcal B_\pi\) when the stated
window and tolerance criteria are met; smaller admissible block
algebras can also exist. It proves that
\(\|\Phi_t^*(Y)-Y\|\le t\,\Gamma_{\mathrm{in}}(\pi)\|Y\|\) for \(Y\in\mathcal B_\pi\), with
\(\Gamma_{\mathrm{in}}\) the block-width bound of HAFF-B~\cite{HAFFB}:
\[
\Gamma_{\mathrm{in}}(\pi)
=\max_{C\in\pi}\frac12\sum_k\gamma_k
 \bigl(\max_{x\in C}a_k(x)-\min_{x\in C}a_k(x)\bigr)^2.
\]
The maximum is taken after summing the squared widths of each block; it is
not generally the largest pairwise decoherence rate within that block.
The pinching
\(E_\pi\) onto \(\mathcal B_\pi\) commutes with the Schur multiplier \(\Phi_t^*\), so
\[
\|E_\pi\circ\Phi_t^*-E_\pi\|=\|\Phi_t^*\circ E_\pi-E_\pi\|\le t\,\Gamma_{\mathrm{in}}(\pi),
\]
and the generator has no matrix elements between \(\mathcal B_\pi\) and its complement.
If the source state lies in \(\mathcal B_\pi\), then \(\omega\circ E_\pi=\omega\) and
\(\mathcal B_\pi\) is invariant under the modular group. On the window, the selected
algebra plays the part of \(\Fix(\mathcal E)\) and the pinching is a Nakajima--Zwanzig
projection that persists up to \(t\Gamma_{\mathrm{in}}\). In this class the memory
kernel between the selected algebra and the rest vanishes exactly; a persistent record
process with genuine memory needs non-commuting couplings or environments that are
revisited, as in Paper~F of HAFF~\cite{HAFFF}.

\section{Structure of the stable hull}\label{sec:hull}

\begin{theorem}[The hull]\label{thm:hull}
Let \(\Gamma\supseteq\Gamma_\sigma\) and let \(S\subseteq\AU\) be a seed. Write
\(\mathcal H=\Hull_\Gamma(S)\).
\begin{enumerate}[label=(\alph*),leftmargin=2em]
\item \(\mathcal H\) is \(\Gamma\)-admissible and is contained in every
  \(\Gamma\)-admissible algebra that contains \(S\). With \(M_0=W^*(S)\) and
  \(M_{n+1}=W^*\bigl(M_n\cup\bigcup_{\gamma\in\Gamma}\gamma(M_n)\bigr)\),
  \(\mathcal H=\bigl(\bigcup_nM_n\bigr)''\).
\item There is a unique normal conditional expectation \(E_{\mathcal H}:\AU\to\mathcal H\)
  with \(\omega\circ E_{\mathcal H}=\omega\); it is faithful. The modular group of
  \(\omega|_{\mathcal H}\) is \(\sigma|_{\mathcal H}\). The projection
  \(e_{\mathcal H}\) onto \(\overline{\mathcal H\Omega}\) commutes with \(\Delta^{it}\)
  and \(J\), and \(\Delta|_{\overline{\mathcal H\Omega}}\),
  \(J|_{\overline{\mathcal H\Omega}}\) are the modular objects of
  \((\mathcal H,\Omega)\) on that subspace.
\item \(S\subseteq S'\) implies \(\Hull_\Gamma(S)\subseteq\Hull_\Gamma(S')\), and
  \(W^*(S)\subseteq\mathcal H\); in particular \(\mathcal H\neq\mathbb C\mathbf 1\) as
  soon as \(S\) contains a nonscalar element.
\end{enumerate}
\end{theorem}

\begin{proof}
(a) An intersection of admissible von Neumann algebras is admissible, and \(\AU\) is
admissible. Every admissible algebra containing \(S\) contains each \(M_n\) by
induction. Each \(\gamma\in\Gamma\) is \(\sigma\)-weakly continuous and maps the
\(*\)-algebra \(\bigcup_nM_n\) into itself, hence maps its \(\sigma\)-weak closure, which
is \((\bigcup_nM_n)''\), into that closure. (b) \(\mathcal H\) is \(\sigma\)-invariant,
so Takesaki's theorem gives \(E_{\mathcal H}\) and the restriction
property~\cite{Takesaki1972}. Since \(\Delta^{it}\mathcal H\Omega=\sigma_t(\mathcal H)\Omega
=\mathcal H\Omega\), \(e_{\mathcal H}\) commutes with \(\Delta^{it}\), hence with every
function of \(\Delta\). For \(x\in\mathcal H\), \(Jx^*\Omega=\Delta^{1/2}x\Omega
=\Delta^{1/2}e_{\mathcal H}x\Omega=e_{\mathcal H}\Delta^{1/2}x\Omega\), so \(J\) maps the
dense subspace \(\mathcal H\Omega\) of \(\overline{\mathcal H\Omega}\) into it. The
restricted operators are the modular objects of \((\mathcal H,\Omega)\) because the
Tomita operator of \(\mathcal H\) is the restriction of that of \(\AU\) to a core that
is invariant under \(\Delta^{it}\). (c) is immediate from (a).
\end{proof}

By part (b) the hull is never an arbitrary subalgebra. It carries its own modular
data, and those data are the source's. When the
hull is later read for geometry, the modular flow that enters is the restriction of
the flow of the source state; nothing is added.

For stable seeds of a stationary source the hull has a simple form.

\begin{theorem}[Stable seeds]\label{thm:stableseed}
Let the source be stationary and let \(S\subseteq\Fix(\mathcal E)\).
All statements concerning \(\Gamma_3\) also assume an ambient positive translation structure and
\(\mathcal E_t\circ\Ad U(a)=\Ad U(a)\circ\mathcal E_t\) for \(a\ge0\).
\begin{enumerate}[label=(\alph*),leftmargin=2em]
\item \(\Hull_{\Gamma_2}(S)=\Hull_{\Gamma_\sigma}(S)
  =W^*\bigl(\bigcup_s\sigma_s(S)\bigr)\subseteq\Fix(\mathcal E)\). If there is an ambient
  Cut with the stated covariance, then
  \(\Hull_{\Gamma_3}(S)\) is the smallest von Neumann algebra containing \(S\) that is
  invariant under \(\sigma\) and \(\Ad U(a)\), \(a\ge0\), and it lies in
  \(\Fix(\mathcal E)\).
\item \(W^*(S)\subseteq\Hull(S)\subseteq\Fix(\mathcal E)\) for \(\Hull=\Hull_{\Gamma_2}\)
  or \(\Hull_{\Gamma_3}\). The hull is nontrivial if \(S\) contains a nonscalar
  element, and proper if \(\Fix(\mathcal E)\neq\AU\).
\item \(\Hull(S\cup S')=\Hull(S)\vee\Hull(S')\) for stable seeds \(S,S'\).
\item \(E_{\mathcal H}\circ\mathcal E_t=\mathcal E_t\circ E_{\mathcal H}=E_{\mathcal H}\)
  for \(\mathcal H=\Hull(S)\): the expectation onto the hull is a persistent projection
  with no back-flow.
\end{enumerate}
\end{theorem}

\begin{proof}
(a) \(\Fix(\mathcal E)\) is a \(\sigma\)-invariant von Neumann algebra
(Theorem~\ref{thm:stationary}) containing \(S\), so it contains the modular hull. The
noise acts as the identity on \(\Fix(\mathcal E)\), so the modular hull is
\(\Gamma_2\)-admissible, and it is contained in every \(\Gamma_2\)-admissible algebra
containing \(S\). If \(\mathcal E_t\) commutes with \(\Ad U(a)\), then
\(U(a)xU(a)^*\) is fixed whenever \(x\) is, so \(\Fix(\mathcal E)\) is invariant under
forward translations and the same argument applies. (b) follows from (a). (c) The
algebra generated by two admissible subalgebras of \(\Fix(\mathcal E)\) is invariant
under the automorphisms \(\sigma_s\) and the normal endomorphisms \(\Ad U(a)\), and the
noise is the identity on it. (d) For \(a\in\mathcal H\subseteq\Fix(\mathcal E)\) and
\(x\in\AU\), using bimodularity, Theorem~\ref{thm:stationary}(a) and stationarity,
\[
\omega\bigl(a\,E_{\mathcal H}(\mathcal E_t(x))\bigr)=\omega(a\,\mathcal E_t(x))
=\omega(\mathcal E_t(ax))=\omega(ax)=\omega\bigl(a\,E_{\mathcal H}(x)\bigr).
\]
Hence \(E_{\mathcal H}(\mathcal E_t(x))-E_{\mathcal H}(x)\in\mathcal H\) is orthogonal to
\(\mathcal H\Omega\) and vanishes because \(\Omega\) is separating. The second identity
holds because the range of \(E_{\mathcal H}\) is fixed.
\end{proof}

Part (c) answers one of the obstructions listed in HAFF-D: a completely positive map
that is not multiplicative can preserve two subalgebras without preserving their
join~\cite{HAFFD}. For stable seeds of a stationary source the noise is multiplicative
on the relevant algebra, and the hull is a closure operator that preserves joins.

\begin{proposition}[Records that commute with the state]\label{prop:records}
Let the source be stationary and let \(S\subseteq\Fix(\mathcal E)\) lie in the
centralizer \(\AU^\omega=\{x:\sigma_s(x)=x\ \forall s\}\). Then
\(\Hull_{\Gamma_2}(S)=W^*(S)\). In the commuting-coupling class of HAFF-B with a
stationary faithful state, the block projections, which generate the record algebra,
lie in the centralizer; the stable hull of any set of record projections is the
abelian algebra they generate, and it contains no proper half-sided modular inclusion.
\end{proposition}

\begin{proof}
The centralizer is a von Neumann algebra on which \(\sigma\) acts trivially, so
\(W^*(S)\) is \(\sigma\)-invariant; it lies in \(\Fix(\mathcal E)\), where the noise is
trivial. In the commuting class a stationary state has vanishing entries between basis
vectors with different eigenvalue tuples of the couplings~\cite{HAFFB}, so it commutes
with the block projections. The modular group of a faithful state on an abelian algebra
is trivial, so the translation generator of a half-sided modular inclusion inside it
vanishes~\cite{AZ} and the inclusion is an equality.
\end{proof}

These centralizer records do not grow under \(\Gamma_2\). The thermal dephasing chain of HAFF-D at zero
transverse field, where the diagonal algebra is fixed and the modular flow acts
trivially on it, is an instance~\cite{HAFFD}. For a stable seed in a stationary source,
\(\Gamma_2\)-growth requires elements moved by the modular flow. Consciousness~I gives the basic example: in
\(M_2\otimes M_2\) with state \(\operatorname{diag}(0.8,0.2)\otimes\frac12\mathbf 1\) and
depolarising noise on the second factor, \(X\otimes\mathbf 1\) is fixed by the noise
and its hull is \(M_2\otimes\mathbf 1\)~\cite{PaperI}.

In finite dimension the growth is bounded by the state.

\begin{proposition}[Finite dimension]\label{prop:finite}
If \(\AU=M_n\), \(\omega=\operatorname{Tr}(\rho\,\cdot)\) is stationary and
\(S\subseteq\Fix(\mathcal E)\), then
\(W^*(S)\subseteq\Hull_{\Gamma_2}(S)\subseteq W^*(S\cup\{\rho\})\cap\Fix(\mathcal E)\).
\end{proposition}

\begin{proof}
\(\sigma_s(x)=\rho^{is}x\rho^{-is}\in W^*(S\cup\{\rho\})\); combine with
Theorem~\ref{thm:stableseed}.
\end{proof}

For three qubits under the collective dephasing \(X_1+X_2+X_3\) of HAFF-B, with a
stationary Gibbs-type state in the block algebra and the in-block seed \(X_1-X_2\), the
seed generates a 3-dimensional algebra, its hull has dimension 10, equal to that of
\(W^*(S\cup\{\rho\})\), and the fixed algebra has dimension 20 (check~4 of
Appendix~\ref{app:checks}). In type~III the picture changes.

\begin{proposition}[Chiral vacuum]\label{prop:chiral}
Let \(\AU=\mathcal A(0,\infty)\) be the half-line algebra of a M\"obius covariant net
with unique vacuum \(\Omega\). Then the centralizer of \(\omega\) is \(\mathbb C\mathbf 1\),
and every \(\sigma\)-invariant von Neumann subalgebra other than \(\mathbb C\mathbf 1\) is
infinite-dimensional.
\end{proposition}

\begin{proof}
By the Bisognano--Wichmann property the modular group is the dilation group
\(D(t)\)~\cite{BGL}, and \(D(t)U(a)D(t)^*=U(e^{-2\pi t}a)\) for the translations. If
\(D(t)\xi=\xi\) for all \(t\), then
\(\langle\xi,U(a)\xi\rangle=\langle\xi,U(e^{-2\pi t}a)\xi\rangle\to\|\xi\|^2\) as
\(t\to\infty\), so \(U(a)\xi=\xi\) and \(\xi\) is a multiple of \(\Omega\). For \(x\) in
the centralizer, \(x\Omega\) is dilation invariant, hence \(x\Omega=c\Omega\) and
\(x=c\mathbf 1\). On the orthogonal complement of \(\Omega\) the representation of the
translation--dilation group has positive translation generator and no invariant
vectors, so it is a multiple of the unique irreducible representation of this kind, in
which the dilation generator has absolutely continuous spectrum. A finite-dimensional
\(\sigma\)-invariant subalgebra \(\mathcal N\) would make \(\mathcal N\Omega\) a
finite-dimensional subspace invariant under \(\Delta^{it}\), hence spanned by
eigenvectors of \(\log\Delta\); these are multiples of \(\Omega\), so
\(\mathcal N=\mathbb C\mathbf 1\).
\end{proof}

In a chiral vacuum every nonscalar seed is moved by the modular flow and generates an
infinite-dimensional hull. Thus in this chiral vacuum the only
centralizer seeds of Proposition~\ref{prop:records} are scalar.
Other type~III source states can have nontrivial centralizers.

\section{Cuts in the hull}\label{sec:cut}

Consciousness~I needs a Cut inside the accessible algebra. A hull closed under
forward translations inherits a half-sided modular inclusion from the ambient positive translation structure;
it is proper when the hull is nonscalar and the invariant vector is unique.

\begin{theorem}[Cut inheritance]\label{thm:cut}
Let the source carry an ambient positive translation structure \(U\), and let \(\mathcal H\subseteq\AU\) be a von
Neumann subalgebra with \(\sigma_s(\mathcal H)=\mathcal H\) for all \(s\) and
\(U(a)\mathcal HU(a)^*\subseteq\mathcal H\) for \(a\ge0\); every \(\Gamma_3\)-hull has
these properties. Put \(\mathcal K=\overline{\mathcal H\Omega}\) and
\(\mathcal N_{\mathcal H}=U(1)\mathcal HU(1)^*\).
\begin{enumerate}[label=(\alph*),leftmargin=2em]
\item \(\mathcal K\) reduces \(U(a)\) for every \(a\in\mathbb R\), as well as
  \(\Delta^{it}\) and \(J\).
\item On \(\mathcal K\), \((\mathcal N_{\mathcal H}\subset\mathcal H,\Omega)\) is a
  half-sided modular inclusion, and its translation group is \(U|_{\mathcal K}\).
\item \(\mathcal N_{\mathcal H}\neq\mathcal H\) if and only if \(U|_{\mathcal K}\neq\mathbf 1\).
  If \(\Omega\) is, up to scalars, the only translation-invariant vector, then
  \(\mathcal N_{\mathcal H}\neq\mathcal H\) whenever \(\mathcal H\neq\mathbb C\mathbf 1\).
\end{enumerate}
In particular, if \(\Omega\) is, up to scalars, the only \(U\)-invariant vector, as in the vacuum representation of a chiral net, then for every seed containing a nonscalar element the \(\Gamma_3\)-hull
contains a Cut in its vacuum-cyclic representation, and its translations are the
restricted ambient translations.
\end{theorem}

\begin{proof}
(a) For \(a\ge0\) and \(x\in\mathcal H\), \(U(a)x\Omega=U(a)xU(a)^*\Omega\in\mathcal H\Omega\).
By Theorem~\ref{thm:hull}(b), \(J\mathcal K=\mathcal K\) and \(\Delta^{it}\mathcal K
=\mathcal K\). Borchers' relation \(U(-a)=JU(a)J\) then gives invariance under
negative translations.

(b) \(\mathcal N_{\mathcal H}\Omega=U(1)\mathcal H\Omega\) is dense in
\(U(1)\mathcal K=\mathcal K\), and \(\Omega\) is separating for every subalgebra of
\(\AU\). The modular group of \((\mathcal H,\Omega)\) on \(\mathcal K\) is
\(\Delta^{it}|_{\mathcal K}\) by Theorem~\ref{thm:hull}(b). For \(t\ge0\), Borchers'
relation \(\Delta^{-it}U(1)\Delta^{it}=U(e^{2\pi t})\) gives
\[
\Delta^{-it}\mathcal N_{\mathcal H}\Delta^{it}
=U(e^{2\pi t})\mathcal HU(e^{2\pi t})^*
=U(1)\bigl[U(e^{2\pi t}-1)\mathcal HU(e^{2\pi t}-1)^*\bigr]U(1)^*
\subseteq\mathcal N_{\mathcal H},
\]
which is the half-sided condition. The modular operator of \(\mathcal N_{\mathcal H}\)
is \(U(1)\Delta U(1)^*\) on \(\mathcal K\), so
\(\Delta_{\mathcal N_{\mathcal H}}^{it}\Delta_{\mathcal H}^{-it}=U(1-e^{-2\pi t})|_{\mathcal K}\).
This is the relation that defines the translation group of the inclusion in the
Wiesbrock--Araki--Zsid\'o theorem~\cite{Wiesbrock1993,AZ}, which is therefore
\(U|_{\mathcal K}\).

(c) If \(\mathcal N_{\mathcal H}=\mathcal H\), applying \(\sigma_t\) gives
\(U(b)\mathcal HU(b)^*=\mathcal H\) for all \(b>0\), hence for all \(b\). On
\(\mathcal K\) the group \(U\) then implements automorphisms of \(\mathcal H\) and of
its commutant \(\mathcal H'\cap\mathcal B(\mathcal K)\), with positive generator and
fixed vector \(\Omega\). Borchers' theorem applied to \(\mathcal H\) and to
\(\mathcal H'\), whose modular operator is \(\Delta_{\mathcal H}^{-1}\), gives both
\(\Delta_{\mathcal H}^{it}U(a)\Delta_{\mathcal H}^{-it}=U(e^{-2\pi t}a)\) and
\(\Delta_{\mathcal H}^{it}U(a)\Delta_{\mathcal H}^{-it}=U(e^{2\pi t}a)\) on
\(\mathcal K\). Hence \(U(a)=U(e^{4\pi t}a)\) for all \(t\), and letting
\(t\to-\infty\) gives \(U|_{\mathcal K}=\mathbf 1\). Conversely, if \(U|_{\mathcal K}=\mathbf 1\)
then \(\mathcal N_{\mathcal H}\) and \(\mathcal H\) act identically on \(\mathcal K\). If
\(\Omega\) is the only invariant vector, \(U|_{\mathcal K}=\mathbf 1\) forces
\(\mathcal K=\mathbb C\Omega\) and \(\mathcal H=\mathbb C\mathbf 1\).
\end{proof}

In Consciousness~I the Virasoro algebra \(B\) of the half-line is fixed by the noise,
\(\Omega\) is not cyclic for \(B\) on the ambient space, and the Cuts used for the
geometry are the subalgebras \(\mathrm{Vir}_1(a,\infty)\) of \(B\) on its own vacuum
space~\cite{PaperI}. Under the unique-invariant-vector hypothesis of
Theorem~\ref{thm:cut}, the \(\Gamma_3\)-hull of any nonscalar seed carries a Cut
on its vacuum-cyclic subspace, and the modular
orbit of that Cut generates the hull by the modular-orbit generation lemma
of Consciousness~I~\cite{PaperI}. If the inclusion is
standard, the Guido--Longo--Wiesbrock correspondence turns it into a strongly additive
conformal net with translations \(U|_{\mathcal K}\) and dilations
\(\Delta^{it}|_{\mathcal K}\)~\cite{GLW}.

Two points limit what this achieves. The translations are the ambient ones; the hull
selects a subspace and an algebra on which they act, and does not produce them.
Standardness is not automatic; it holds in the symmetric examples of the next section.
Since these translation-stable hulls carry the same translation and dilation groups, restricted to
different subspaces, the translation--dilation group alone cannot distinguish the
geometries of different seeds. An invariant of the whole triple is needed, as
Consciousness~I already notes~\cite{PaperI}.

\section{Seeds and symmetry}\label{sec:symmetry}

This section answers three questions about the hull: whether it is unique relative to
the seed, how it changes with the seed, and whether different seeds give different
geometries. The answer is governed by symmetry.

\subsection{A Galois law for hulls}

\begin{lemma}[Vacuum-space criterion]\label{lem:vacspace}
Let \(\mathcal B\subseteq\mathcal C\subseteq\AU\) be \(\sigma\)-invariant von Neumann
algebras. Then \(\mathcal B=\mathcal C\) if and only if
\(\overline{\mathcal B\Omega}=\overline{\mathcal C\Omega}\).
\end{lemma}

\begin{proof}
Let \(E_{\mathcal B}\) be the \(\omega\)-preserving expectation onto \(\mathcal B\). For
\(x\in\AU\) and \(y\in\mathcal B\),
\(\langle y\Omega,(x-E_{\mathcal B}(x))\Omega\rangle=\omega(y^*x)-\omega(E_{\mathcal B}(y^*x))=0\),
so \(e_{\mathcal B}x\Omega=E_{\mathcal B}(x)\Omega\). If the two subspaces agree and
\(x\in\mathcal C\), then \(x\Omega=E_{\mathcal B}(x)\Omega\) and \(x=E_{\mathcal B}(x)\in\mathcal B\).
\end{proof}

\paragraph{Symmetric sources.} Let \(\AU\) be a factor with separable predual and \(G\)
a compact group with a strongly continuous unitary representation \(u\) on
\(\mathcal H_U\) such that \(u(g)\Omega=\Omega\) and
\(\alpha_g=\Ad u(g)\) maps \(\AU\) onto itself; if an ambient positive translation structure is present, \(u(g)\)
commutes with \(U(a)\). Then \(u(g)\) commutes with the Tomita operator, hence with
\(\Delta\), and \(\alpha_g\circ\sigma_s=\sigma_s\circ\alpha_g\). The action is
\emph{minimal} if \(g\mapsto\alpha_g\) is injective and \((\AU^G)'\cap\AU=\mathbb C\mathbf 1\). The \emph{stabilizer} of a seed
is the closed subgroup
\(\Stab(S)=\{g\in G:\alpha_g(x)=x\ \text{for all}\ x\in S\}\).
Izumi, Longo and Popa proved that for a minimal action of a compact group on a factor
with separable predual, \(K\mapsto\AU^K\) is a bijection from the closed subgroups of
\(G\) onto the von Neumann algebras between \(\AU^G\) and \(\AU\)~\cite{ILP1998}.

\begin{theorem}[Symmetry fixes the hull]\label{thm:galois}
Let \(G\) act minimally as above, let \(\Gamma\supseteq\Gamma_\sigma\), and let \(S\)
be a seed such that
\begin{enumerate}[label=(\roman*),leftmargin=2em]
\item every \(\gamma\in\Gamma\) maps \(\AU^{\Stab(S)}\) into itself, and
\item \(\Hull_\Gamma(S)\supseteq\AU^G\).
\end{enumerate}
Then \(\Hull_\Gamma(S)=\AU^{\Stab(S)}\). Two seeds satisfying (i) and (ii) have the same
hull if and only if they have the same stabilizer. Condition~(i) holds if every
\(\gamma\in\Gamma\) commutes with every \(\alpha_g\); for a stable seed of a stationary
source it also holds when \(\Gamma\) consists of \(\sigma\), forward translations and a
noise with \(\Fix(\mathcal E)=\AU^{H_0}\) for a closed subgroup \(H_0\).
\end{theorem}

\begin{proof}
By (i), \(\AU^{\Stab(S)}\) is \(\Gamma\)-admissible, and it contains \(S\); hence
\(\Hull_\Gamma(S)\subseteq\AU^{\Stab(S)}\). By (ii) and the Izumi--Longo--Popa theorem,
\(\Hull_\Gamma(S)=\AU^{K'}\) for a closed subgroup \(K'\). Since
\(S\subseteq\AU^{K'}\), \(K'\subseteq\Stab(S)\), so
\(\AU^{\Stab(S)}\subseteq\AU^{K'}=\Hull_\Gamma(S)\). The second statement is the
injectivity of \(K\mapsto\AU^K\). For the last statement, if \(S\subseteq\AU^{H_0}\)
then \(\Stab(S)\supseteq H_0\), so \(\AU^{\Stab(S)}\subseteq\Fix(\mathcal E)\), where the
noise is the identity; \(\sigma\) and the translations commute with \(\alpha_g\).
\end{proof}

The hull is therefore unique relative to the seed in a strong sense: it depends on the
seed only through the symmetry the seed leaves unbroken. Adding observables to a seed
can only shrink the stabilizer and enlarge the hull, and observables that do not change
the stabilizer add nothing.

The law has no counterpart in type~I. Every automorphism of a
type~I factor is inner, so each \(\alpha_g\) is implemented by a unitary \(v(g)\in\AU\); these unitaries
commute with \(\AU^G\), so a nontrivial action has
\((\AU^G)'\cap\AU\supseteq\{v(g)\}''\neq\mathbb C\mathbf 1\), and no nontrivial action is
minimal. Distinct stabilizers can then give the same algebra: in
\(M_2\) with the adjoint action of \(SO(3)\), both \(O(2)\) and \(SO(3)\) fix only the
scalars (check~B.1.6).

\subsection{The SU(2) level-one current net}

Let \(\mathcal A\) be the SU(2)\(_1\) current net, \(\AU=M=\mathcal A(0,\infty)\) with the
vacuum, the translations \(U\) as ambient positive translation structure and the dilations as modular
group~\cite{BGL}. The group \(SO(3)\) acts by the adjoint action, as in Assumption~4.1 of
Consciousness~I~\cite{PaperI}. Write \(J_n(f)=\sum_an_aJ^a(f)\) for a unit vector \(n\)
and a real test function \(f\), so that \(\alpha_g(J_n(f))=J_{R(g)n}(f)\), and
\(\mathcal A_{J_n}(I)=\{e^{iJ_n(h)}:\operatorname{supp}h\subset I\}''\) for the U(1) current
subnet of \(J_n\). We use four known facts.
\begin{enumerate}[label=(F\arabic*),leftmargin=2.5em]
\item \(M^{SO(3)}=\Vir_1(0,\infty)\), the half-line algebra of the Virasoro net of the
  Sugawara stress tensor (Rehren~\cite{Rehren}).
\item The Sugawara stress tensor equals the stress tensor of each current
  \(J_n\): their difference commutes with \(J_n\) and is a positive-energy Virasoro field
  of central charge \(1-1=0\), hence zero~\cite{GKO1986}. Thus
  \(\Vir_1(I)\subseteq\mathcal A_{J_n}(I)\)~\cite{Carpi1999}.
\item On \(\overline{\mathcal A_{J_n}(0,\infty)\Omega}\) the field \(J_n\) is a free
  current: the vacuum is quasi-free on \(\mathcal A_{J_n}\), with one-particle space the
  completion of the real test functions in the norm
  \(\|h\|^2=c\int_0^\infty k\,|\hat h(k)|^2\,dk\), and Weyl relations
  \(e^{iJ_n(h_1)}e^{iJ_n(h_2)}=e^{i\theta(h_1,h_2)}e^{iJ_n(h_1+h_2)}\) with real
  \(\theta\) hold on the whole space.
\item \(\Vir_1\) and \(\mathcal A\) are strongly additive~\cite{Carpi,Xu}, and the
  fixed-point subnet of a strongly additive net under a compact group is strongly
  additive~\cite{Xu}.
\end{enumerate}

\begin{proposition}[Minimality]\label{prop:minimal}
\(SO(3)\) acts minimally on \(M\): the action is faithful and \(\Vir_1(0,\infty)'\cap M=\mathbb C\mathbf 1\). This is known~\cite{Carpi1999,Xu}; we include a proof.
\end{proposition}

\begin{proof}
Let \(x\in M\) commute with \(\Vir_1(0,\infty)\). For an interval \(I\) inside the
half-line, \(\Vir_1(I)\subseteq\Vir_1(0,\infty)\) commutes with \(x\); for \(I\) in the
complement, \(\Vir_1(I)\subseteq\mathcal A(I)\subseteq M'\) by locality. An interval
that contains \(0\) or \(\infty\) splits, after removing these points, into intervals of
the first two kinds, and by strong additivity of \(\Vir_1\) its algebra is generated by
theirs. So \(x\) commutes with \(\Vir_1(I)\) for every interval, hence with
\(V(\varphi)\) for every diffeomorphism supported in an interval, hence, since such
diffeomorphisms generate \(\mathrm{Diff}_+(S^1)\), with the M\"obius group. Then
\(x\Omega\) is M\"obius invariant, \(x\Omega=c\Omega\) and \(x=c\mathbf 1\). Faithfulness: \(\alpha_g(J_n(f))=J_{R(g)n}(f)\) and the vectors \(J^a(f)\Omega\), \(a=1,2,3\), are orthogonal and nonzero, so \(\alpha_g=\mathrm{id}\) forces \(R(g)=1\).
\end{proof}

\begin{lemma}[A current seed]\label{lem:current}
For every unit vector \(n\) and every nonzero real
\(g\in C_c^\infty((0,\infty))\), the modular hull of
\(O_{n,g}=\tfrac2\pi\arctan J_n(g)\) is \(\mathcal A_{J_n}(0,\infty)\).
\end{lemma}

\begin{proof}
\(\mathcal A_{J_n}(0,\infty)\) is invariant under dilations and contains \(O_{n,g}\), so
it contains the modular hull. Conversely, the hull contains \(e^{itJ_n(g)}\) for all
real \(t\), and dilations act by \(e^{iJ_n(h)}\mapsto e^{iJ_n(h_\lambda)}\),
\(h_\lambda(x)=h(x/\lambda)\). By the Weyl relations the hull contains
\(e^{iJ_n(h)}\) for every \(h\) in the real span \(L\) of the dilates of \(g\); let
\(\mathcal B_L\) be the von Neumann algebra these generate. It is dilation invariant and
lies in \(\mathcal A_{J_n}(0,\infty)\).

\(L+iL\) is dense in the one-particle space. With \(k=e^u\) and
\(\psi(u)=e^u\hat g(e^u)\), the one-particle norm is the \(L^2(\mathbb R,du)\) norm and a
dilation is a translation of \(\psi\). The Fourier transform
\(\tilde\psi(s)=\int_0^\infty k^{-is}\hat g(k)\,dk\) is the restriction to
\(\operatorname{Re}z=1\) of the Mellin transform \(\int_0^\infty k^{z-1}\hat g(k)\,dk\),
which is analytic for \(\operatorname{Re}z>0\) because \(\hat g\) is bounded and
rapidly decreasing, and is not identically zero because \(g\neq0\). Its zeros on the
line are isolated, so \(\tilde\psi\neq0\) almost everywhere, and by Wiener's theorem
the translates of \(\psi\) span a dense subspace~\cite{Wiener1932}.

The vectors \(e^{iJ_n(h)}\Omega\), \(h\in L\), are multiples of the coherent vectors
\(\exp(ih)\). Differentiating \(\exp(i\sum_jt_jh_j)\) in the real parameters \(t_j\) at
zero produces all symmetric tensor products of elements of \(L\), whose complex span is
dense in each \(n\)-particle space because \(L+iL\) is dense. Hence
\(\overline{\mathcal B_L\Omega}\) is the Fock space
\(\overline{\mathcal A_{J_n}(0,\infty)\Omega}\), and Lemma~\ref{lem:vacspace} gives
\(\mathcal B_L=\mathcal A_{J_n}(0,\infty)\).
\end{proof}

\begin{lemma}[A stress seed under modular evolution]\label{lem:stressmod}
Let \(0<a<b\), \(\alpha=\log a\), \(\beta=\log b\), and
\(g\in C_c^\infty(\mathbb R,\mathbb R)\) have support \([\alpha,\beta]\).
Suppose \(g>0\) on \((\alpha,\beta)\), is strictly increasing near \(\alpha\),
and strictly decreasing near \(\beta\). For
\(f(x)=xg(\log x)\) and \(O_f=\tfrac2\pi\arctan T(f)\),
\[
\Hull_{\Gamma_\sigma}(\{O_f\})=\Vir_1(0,\infty).
\]
The one-sided orbit \(N^-_{O_f}=W^*(\sigma_s(O_f):s\le0)\), restricted to
the Virasoro vacuum space, is a proper standard Cut and satisfies
\(\Vir_1(b,\infty)\subseteq N^-_{O_f}\subseteq\Vir_1(a,\infty)\).
\end{lemma}

\begin{proof}
These are the single-orbit Cut theorem and full modular-hull corollary of
Consciousness~I~\cite{PaperI}. The ambient modular group restricts to the
Virasoro vacuum modular group by (F1). Every stress flow used in that
construction belongs to the ambient algebra; its linear combinations,
conjugations and fixed-compact-support smooth limits retain their strong
closure there. The construction therefore puts
\(\Vir_1(b,\infty)\) in the one-sided orbit algebra without an independent
forward-translation permission. Modular invariance of the full hull then
puts \(\Vir_1(\lambda b,\infty)\) in it for every \(\lambda>0\). These
algebras generate \(\Vir_1(0,\infty)\), which is also an upper bound
because it contains the seed and is dilation invariant.
Standardness and properness refer to the Virasoro vacuum space, not the
larger SU(2)\(_1\) Hilbert space.
\end{proof}

Thus a current seed always has modular-only generation, while a stress
seed has it under the stated bump conditions. For an arbitrary nonzero
real stress smearing, the general stress-generation theorem of
Consciousness~I still uses dilations and forward translations.
None of these closure results supplies a physical stress-flow gate
from an unsharp effect.

\begin{theorem}[Hulls in the SU(2)\(_1\) net]\label{thm:su2}
Let \(S\) be a seed containing at least one of the following:
\begin{enumerate}[label=(\roman*),leftmargin=2em]
\item a stress observable \(O_f=\frac2\pi\arctan T(f)\) with arbitrary
  \(0\neq f\in C_c^\infty((0,\infty),\mathbb R)\), with \(\Gamma=\Gamma_3\);
\item a stress observable \(O_f\) satisfying Lemma~\ref{lem:stressmod},
  with \(\Gamma\in\{\Gamma_\sigma,\Gamma_2,\Gamma_3\}\);
\item a current observable \(O_{n,g}\), with
  \(\Gamma\in\{\Gamma_\sigma,\Gamma_2,\Gamma_3\}\).
\end{enumerate}
Let the noise be absent, the Haar noise
\(D^{SO(3)}\), or, for a stable seed, a Haar noise \(D^{H_0}\) of Consciousness~I.
Put \(K=\Stab(S)\). Then:
\begin{enumerate}[label=(\alph*),leftmargin=2em]
\item \(\Hull_\Gamma(S)=M^K\);
\item \(\overline{M^K\Omega}=\mathcal H_U^K\), the \(K\)-invariant vectors, and its
  character is
  \[
  \operatorname{Tr}_{\mathcal H_U^K}q^{L_0}
  =\frac{1}{\phi(q)}\sum_{j\ge0}m_j(K)\bigl(q^{j^2}-q^{(j+1)^2}\bigr),
  \qquad m_j(K)=\dim V_j^K,
  \]
  where \(V_j\) is the spin-\(j\) representation of \(SO(3)\) and
  \(\phi(q)=\prod_{n\ge1}(1-q^n)\);
\item \(U(1)M^KU(1)^*\subset M^K\) is a standard Cut on \(\mathcal H_U^K\), and the net
  it reconstructs is the fixed-point subnet \(I\mapsto\mathcal A(I)^K\) with the
  restricted M\"obius representation.
\end{enumerate}
\end{theorem}

\begin{proof}
(a) Every \(\Gamma\)-admissible algebra is invariant under \(\sigma\);
a \(\Gamma_3\)-admissible algebra is also stable under forward translations.
In case (i), the general stress-generation theorem of
Consciousness~I~\cite{PaperI} puts \(\Vir_1(0,\infty)\) in the hull.
In case (ii), Lemma~\ref{lem:stressmod} gives the same inclusion using
only modular evolution. In case (iii), it follows from
Lemma~\ref{lem:current} and (F2). By (F1) this is \(M^{SO(3)}\).
The modular group and the translations commute with \(\alpha_g\); the Haar noise
\(D^{SO(3)}\) commutes with every \(\alpha_g\) by invariance of Haar measure; a noise
\(D^{H_0}\) has fixed algebra \(M^{H_0}\)~\cite{PaperI}. Proposition~\ref{prop:minimal}
and Theorem~\ref{thm:galois} give the claim.

(b) Since \(u(k)\Omega=\Omega\), the projection \(P_K=\int_Ku(k)\,dk\) onto
\(\mathcal H_U^K\) maps \(x\Omega\) to \(Q_K(x)\Omega\in M^K\Omega\), where \(Q_K\) is the
Haar average; as \(M\Omega\) is dense, \(M^K\Omega\) is dense in \(\mathcal H_U^K\). The
vacuum module decomposes under \(SO(3)\times\Vir\) as
\(\bigoplus_{j\ge0}V_j\otimes L(1,j^2)\), and the character of the irreducible
\(c=1\) module \(L(1,j^2)\) is \(q^{j^2}(1-q^{2j+1})/\phi(q)\)~\cite{Rehren,Segal,KR};
taking \(K\)-invariants gives the formula.

(c) \(\mathcal A^K\) is a M\"obius covariant subnet, strongly additive by (F4), and
\(\Omega\) is cyclic on \(\mathcal H_U^K\) for each \(\mathcal A(I)^K\) by the
Reeh--Schlieder property of the subnet in its vacuum representation. Since
\(M^K=\mathcal A^K(0,\infty)\) and \(U(1)M^KU(1)^*=\mathcal A^K(1,\infty)\), the relative
commutant contains \(\mathcal A^K(0,1)\), so the Cut is standard, and the
Guido--Longo--Wiesbrock reconstruction returns \(\mathcal A^K\)~\cite{GLW}. The M\"obius
representation of a M\"obius covariant net is fixed by the modular groups of its
interval algebras~\cite{BGL}, so it is the restriction of the ambient one.
\end{proof}

\begin{theorem}[Seeds are distinguished by geometry]\label{thm:distinguish}
For two seeds \(S_1,S_2\) as in Theorem~\ref{thm:su2}, the Cut triples
\((U(1)\mathcal H_iU(1)^*\subset\mathcal H_i,\Omega)\) on
\(\overline{\mathcal H_i\Omega}\), \(\mathcal H_i=\Hull_\Gamma(S_i)\), are isomorphic if
and only if \(\Stab(S_1)\) and \(\Stab(S_2)\) are conjugate in \(SO(3)\).
\end{theorem}

\begin{proof}
If \(\Stab(S_2)=g\Stab(S_1)g^{-1}\), then \(u(g)\) maps \(M^{\Stab(S_1)}\) onto
\(M^{\Stab(S_2)}\), commutes with \(U\) and \(\Delta\), and fixes \(\Omega\). Conversely,
an isomorphism of standard triples induces an isomorphism of the reconstructed nets
intertwining their M\"obius representations~\cite{GLW,BGL}, so the two characters of
Theorem~\ref{thm:su2}(b) agree. Multiplying by \(\phi(q)\) gives
\(\sum_j(m_j-m_{j-1})q^{j^2}\), from which every \(m_j(K)\) is read off. With
\(\chi_j(\theta)=\sum_{m=-j}^je^{im\theta}\) the character of \(V_j\) at rotation angle
\(\theta\), \(m_j-m_{j-1}=\int(\chi_j-\chi_{j-1})\,d\mu_K=2\int\cos(j\theta)\,d\mu_K(\theta)\)
for \(j\ge1\), where \(\mu_K\) is the image of Haar measure of \(K\) under the rotation
angle \(\theta\in[0,\pi]\). The cosine moments determine \(\mu_K\), and \(\mu_K\)
determines the conjugacy class of \(K\) (Lemma~\ref{lem:angles}).
\end{proof}

The classified triples use the ambient-translation inclusion
$U(1)M^KU(1)^*\subset M^K$. The seed's one-sided orbit
$N^-_{O_f}\subset\mathrm{Vir}_1(0,\infty)$ in
Lemma~\ref{lem:stressmod} is a different standard inclusion.
Its reconstructed translations need not equal $U$.
Consciousness~I's Conjecture T1(b) asserts that this one-sided orbit
is not an ordinary translated half-line; that noncoincidence remains
open. The character invariant here is
proved for the ambient-translation triples.

The identification compares the entire standard Cut triples and their conformal subnets
on the same ambient light ray. Their nonisomorphism does not assert that their underlying
spacetime metrics differ.
The whole character is needed: the multiplicities of \(C_n\) are
\(2\lfloor j/n\rfloor+1\), so its hull and the U(1) hull have the same character below
level \(n^2\). A reader resolving finitely many conformal levels distinguishes only
stabilizers whose characters already differ below that level.

Table~\ref{tab:chars} lists the first coefficients. Some symmetries show up late: the
hulls of seeds with octahedral and icosahedral stabilizers agree with the Virasoro hull
up to level \(15\), and the icosahedral hull differs from the Virasoro hull first at
level \(36\), where the first invariant harmonic of degree six appears.

\begin{table}[ht]
\centering
\small
\begin{tabular}{lll}
\toprule
stabilizer \(K\) & \(m_j(K)\) & character coefficients, levels \(0\)--\(12\)\\
\midrule
\(SO(3)\) & \(\delta_{j0}\) & 1, 0, 1, 1, 2, 2, 4, 4, 7, 8, 12, 14, 21\\
\(O(2)\) & \(1\) for even \(j\), else \(0\) & 1, 0, 1, 1, 3, 3, 6, 7, 12, 14, 22, 27, 40\\
\(SO(2)\) & \(1\) & 1, 1, 2, 3, 5, 7, 11, 15, 22, 30, 42, 56, 77\\
\(D_2\) & & 1, 0, 1, 1, 4, 4, 8, 10, 17, 21, 33, 42, 62\\
\(T\) & & 1, 0, 1, 1, 2, 2, 4, 4, 7, 9, 13, 16, 24\\
\(C_2\) & \(2\lfloor j/2\rfloor+1\) & 1, 1, 2, 3, 7, 9, 15, 21, 32, 44, 64, 86, 121\\
\(\{1\}\) & \(2j+1\) & 1, 3, 4, 7, 13, 19, 29, 43, 62, 90, 126, 174, 239\\
\bottomrule
\end{tabular}
\caption{Vacuum characters of the hulls \(M^K\) in the SU(2)\(_1\) net. The rows for
\(SO(3)\), \(SO(2)\) and \(\{1\}\) are the Virasoro, U(1) current and SU(2)\(_1\) vacuum
characters; the \(O(2)\) row counts partitions into an even number of parts. Computed
in check~B.2 of Appendix~\ref{app:checks}.}
\label{tab:chars}
\end{table}

\begin{corollary}[Dephasing about one axis]\label{cor:oneaxis}
Let the noise be \(D^{H_z}\), with \(H_z\) the rotations about \(e_3\).
The seed and operation family in the hull statements below satisfy
Theorem~\ref{thm:su2}: arbitrary stress smearings use \(\Gamma_3\);
the stated log-bump stress seeds and current seeds also admit
\(\Gamma_\sigma\) and \(\Gamma_2\).
\begin{enumerate}[label=(\alph*),leftmargin=2em]
\item \(\Fix(D^{H_z})=M^{H_z}=\mathcal A_{J_{e_3}}(0,\infty)\), the U(1) current algebra
  of the half-line.
\item A current observable \(O_{n,g}\) is stable if and only if \(n=\pm e_3\). Every
  stable seed that contains a stress or a current observable has one of exactly three
  hulls: \(\mathcal A_{J_{e_3}}(0,\infty)\), its \(\mathbb Z_2\)-orbifold
  \(M^{O(2)_z}\) under \(J_{e_3}\mapsto-J_{e_3}\), or \(\Vir_1(0,\infty)\). They are
  attained by \(\{O_{e_3,g}\}\), by \(\{O_f,x\}\) with
  \(x\in M^{O(2)_z}\setminus\Vir_1(0,\infty)\), and by \(\{O_f\}\), and their characters
  differ at levels \(1\) and \(4\).
\item Under the Haar noise \(D^{SO(3)}\) every stable seed containing a stress
  observable has the hull \(\Vir_1(0,\infty)\). Without noise every conjugacy class of
  closed subgroups occurs as a stabilizer, and two current observables along
  non-parallel axes generate all of \(M\) under the modular flow alone.
\end{enumerate}
\end{corollary}

\begin{proof}
(a) \(\Fix(D^{H_z})=M^{H_z}\)~\cite{PaperI}. Since \(J_n(g)\Omega=\sum_an_aJ^a(g)\Omega\)
with orthogonal nonzero vectors \(J^a(g)\Omega\), \(\alpha_g\) fixes \(O_{e_3,g}\) exactly
when \(R(g)e_3=e_3\); so \(\Stab(O_{e_3,g})=H_z\), and Lemma~\ref{lem:current} with
Theorem~\ref{thm:su2} gives \(\mathcal A_{J_{e_3}}(0,\infty)=M^{H_z}\). (b) The same
argument shows that \(O_{n,g}\in M^{H_z}\) only for \(n=\pm e_3\). A stable seed has
\(\Stab(S)\supseteq H_z\). The closed subgroups of \(SO(3)\) containing \(H_z\) are
\(H_z\), its normalizer \(O(2)_z\) and \(SO(3)\): the identity component of such a
subgroup is \(H_z\) or \(SO(3)\), and in the first case the subgroup normalizes \(H_z\).
The rotation by \(\pi\) about \(e_1\) acts on \(M^{H_z}\) by \(J_{e_3}\mapsto-J_{e_3}\), so
\(M^{O(2)_z}\) is the orbifold. An element \(x\in M^{O(2)_z}\) outside
\(\Vir_1(0,\infty)\) exists by the Galois correspondence, and \(\Stab(\{O_f,x\})\) contains
\(O(2)_z\) and is not \(SO(3)\), so it equals \(O(2)_z\). The characters follow from
\(m_j(H_z)=1\), \(m_j(O(2)_z)=1\) for even \(j\) and \(0\) otherwise, and
\(m_j(SO(3))=\delta_{j0}\). (c) The Haar noise fixes exactly \(M^{SO(3)}\). Without
noise, the self-adjoint part of \(M^K\) is a seed with stabilizer \(K\) by the Galois correspondence, and
\(\Stab(\{O_{n,g},O_{n',g'}\})\) is trivial for non-parallel \(n,n'\).
\end{proof}

Under axis dephasing, stable seeds in this stress/current-containing
class therefore have exactly three ideal hulls, and they are
told apart by the spectrum of the conformal Hamiltonian at the first and fourth levels.
A reader whose ideal seed contains a bounded injective function of
\(J_{e_3}(g)\) has the whole fixed algebra as its hull; a reader that
registers only quantities insensitive to the sign of \(J_{e_3}\), together with the
energy density, gets the orbifold; a reader that registers only energy gets the
Virasoro algebra.

\paragraph{Lifetime of a hull.} Real noise breaks a symmetry slightly. Let
\(\mathcal E_t=D^{H_z}_t\circ D^{SO(3)}_{st}\), with rates \(\kappa\) and \(s\kappa\),
\(0<s\ll1\). The two semigroups commute, because
\(Q_{SO(3)}Q_{H_z}=Q_{H_z}Q_{SO(3)}=Q_{SO(3)}\), and the fixed algebra of the product is
\(\Vir_1(0,\infty)\)~\cite{PaperI}.

\begin{proposition}[A plateau of length \(1/s\)]\label{prop:plateau}
For every closed \(K\supseteq H_z\) with \(M^K\neq\Vir_1(0,\infty)\),
\[
\sup_{x\in M^K,\ \|x\|\le1}\|\mathcal E_t(x)-x\|=c_K\bigl(1-e^{-s\kappa t}\bigr),
\qquad 1\le c_K\le2 .
\]
Thus \(M^K\) is \(\varepsilon\)-stable on \([0,T]\) in the sense of HAFF-B exactly when
\(c_K(1-e^{-s\kappa T})\le\varepsilon\), which for small \(\varepsilon\) is a window of
length \(\varepsilon/(c_Ks\kappa)\).
\end{proposition}

\begin{proof}
On \(M^{H_z}\) the first factor is the identity and
\(D^{SO(3)}_{st}(x)=e^{-s\kappa t}x+(1-e^{-s\kappa t})Q_{SO(3)}(x)\), so
\(\mathcal E_t(x)-x=-(1-e^{-s\kappa t})(x-Q_{SO(3)}(x))\). Put
\(c_K=\sup\{\|x-Q_{SO(3)}(x)\|:x\in M^K,\|x\|\le1\}\le2\). For \(y\in M^K\setminus\Vir_1(0,\infty)\)
the element \(x=(y-Q_{SO(3)}y)/\|y-Q_{SO(3)}y\|\) has \(Q_{SO(3)}(x)=0\), so
\(c_K\ge1\).
\end{proof}

The U(1) hull of a current reader thus lives for a time inversely proportional to the
strength of the symmetry-breaking noise, after which only the Virasoro algebra remains
stable. This is the type~III counterpart of the plateau that HAFF-B computes for the
GHZ geometry~\cite{HAFFB}. A possible maintenance mechanism would correct the symmetry-breaking jumps on their
characteristic timescale \((s\kappa)^{-1}\). The plateau formula establishes the passive
lifetime; it supplies neither a universal correction-rate bound nor the energy cost of
maintaining this type~III hull. Q-RAIF~B gives a field-coupled correcting actuator in a
finite model~\cite{QRAIFB}; extending it to the present hull is a concrete test of HAFF-D's
conjecture that maintained access carries a gravitating stress-energy flux~\cite{HAFFD}.

\section{Ideal seeds and finite reading budgets}\label{sec:foreground}

The symmetry classification concerns ideal operator seeds. A finite
processing budget first restricts a physical reader's instrument menu;
it does not assign a cost to the spectral closure of a current field.
Appendix~\ref{app:budget} gives a separate finite-coordinate coding model.
No dictionary from its coordinate projections to the current or stress
observables is supplied here.

\begin{corollary}[A single ideal current seed]\label{cor:onechannel}
Under the hypotheses of Corollary~\ref{cor:oneaxis}, the ideal spectral
seed $\{O_{e_3,g}\}$, $g\ne0$, has the whole fixed algebra $M^{H_z}$
as its hull. No stable seed has a larger hull. Without noise, two
ideal seeds containing bounded injective functions of nonparallel
current components have $M$ as their hull.
\end{corollary}
\begin{proof}
Lemma~\ref{lem:current} and Theorem~\ref{thm:su2} identify the one-axis
hull. Theorem~\ref{thm:stableseed} bounds every stable hull by
the fixed algebra. Two nonparallel current components have trivial
common stabilizer, so Theorem~\ref{thm:su2}(a) gives $M$.
\end{proof}

This corollary uses ideal bounded functional calculus, all modular times
and von Neumann closure. Proposition~\ref{prop:readbridge} instead
prices an actual finite pointer instrument. Neither statement converts
a one-coordinate processing allowance into permission to implement the
whole current algebra. When a physical menu implements a specified
spectral seed, its ideal hull depends on the stabilizer of that seed;
the menu and its finite cost remain separate inputs.

\section{How a hull responds to excitations}\label{sec:response}

HAFF-D connects a selected algebra to gravity through the response of its modular data
to excitations of the state~\cite{HAFFD}, and Consciousness~I computes the relative
entropy lost when a Cut is displaced~\cite{PaperI}. Whether these responses depend on
the seed is decided by the following facts. For a normal state \(\psi\) and a von
Neumann subalgebra \(\mathcal N\), write \(S_{\mathcal N}(\psi\|\omega)\) for Araki's
relative entropy of the restrictions, and \(\omega_W=\omega(W^*\cdot\,W)\) for a
unitary \(W\).

\begin{proposition}[Response and the hull]\label{prop:response}
Let \(\mathcal H\subseteq\mathcal H'\) be \(\sigma\)-invariant von Neumann subalgebras.
\begin{enumerate}[label=(\alph*),leftmargin=2em]
\item For a unitary \(W\in\mathcal H\),
  \(S_{\mathcal H}(\omega_W\|\omega)=S_{\AU}(\omega_W\|\omega)\), and the same holds with
  the arguments exchanged.
\item For every normal state \(\psi\),
  \(S_{\mathcal H}(\psi\|\omega)\le S_{\mathcal H'}(\psi\|\omega)\). If the right side is
  finite, equality holds if and only if
  \(\psi|_{\mathcal H'}=\psi|_{\mathcal H}\circ E_{\mathcal H}|_{\mathcal H'}\).
\end{enumerate}
\end{proposition}

\begin{proof}
(a) Since \(W\in\mathcal H\) and \(E_{\mathcal H}\) is \(\mathcal H\)-bimodular and
\(\omega\)-preserving, \(\omega_W\circ E_{\mathcal H}=\omega_W\) and
\(\omega\circ E_{\mathcal H}=\omega\). Relative entropy is monotone under restriction to
\(\mathcal H\) and under the channel \(\varphi\mapsto\varphi\circ E_{\mathcal H}\), which
maps the two restrictions back to the two states, so the two values coincide.
(b) The inequality is monotonicity under restriction. The equality condition is
Petz's characterization of sufficient subalgebras~\cite{Petz1986,JencovaPetz2006}: with
\(\omega\circ E_{\mathcal H}=\omega\), the recovery map for \(\omega\) is
\(E_{\mathcal H}\) itself.
\end{proof}

\begin{corollary}[Stress excitations do not see the seed; current excitations do]\label{cor:response}
In the SU(2)\(_1\) net:
\begin{enumerate}[label=(\alph*),leftmargin=2em]
\item for \(W=e^{iT(f)}\) with \(f\in C_c^\infty((0,\infty),\mathbb R)\), every hull \(M^K\) gives the same relative entropy, so the
  coherent-state formula of Hollands and the Cut-displacement theorem of
  Consciousness~I~\cite{Hollands,PaperI} hold verbatim in every hull of
  Theorem~\ref{thm:su2};
\item for \(W=e^{iJ_{e_3}(h)}\) with \(0\neq h\in C_c^\infty((0,\infty),\mathbb R)\),
  \[
  S_{\Vir_1(0,\infty)}(\omega_W\|\omega)
  <S_{\mathcal A_{J_{e_3}}(0,\infty)}(\omega_W\|\omega)
  =2\pi\langle W\Omega,KW\Omega\rangle ,
  \]
  where \(K=-\frac1{2\pi}\log\Delta\) is the dilation generator.
\end{enumerate}
\end{corollary}

\begin{proof}
(a) \(e^{iT(f)}\in\Vir_1(0,\infty)\subseteq M^K\); apply
Proposition~\ref{prop:response}(a). (b) \(W\in\mathcal A_{J_{e_3}}(0,\infty)\), so the
middle term equals the ambient value, which for a unitary of \(M\) is
\(-\langle W\Omega,\log\Delta\,W\Omega\rangle\); it is finite because the coherent state
has finite dilation energy. On \(\mathcal H'=M^{H_z}\) the expectation onto
\(\Vir_1(0,\infty)\) is the Haar average \(Q=Q_{SO(3)}\). For a real
\(k\in C_c^\infty((0,\infty))\) and the odd function \(\arctan\), the rotation by \(\pi\)
about \(e_1\) maps \(\arctan J_{e_3}(k)\) to its negative and leaves \(Q\) unchanged,
so \(Q(\arctan J_{e_3}(k))=0\). In the state \(\omega_W\), \(J_{e_3}(k)\) has the vacuum
law, a centred Gaussian, shifted by a real number \(\mu\) proportional to
\(\int hk'\,dx\), because the commutator of the current is a multiple of the identity;
\(\mu\neq0\) for \(k=h'\). Since \(\arctan\) is odd and strictly
increasing, \(\omega_W(\arctan J_{e_3}(k))\neq0=\omega_W(Q(\arctan J_{e_3}(k)))\). By
Proposition~\ref{prop:response}(b) the inequality is strict.
\end{proof}

For a stress excitation lying in every hull, its relative entropy and modular-energy
response agree across the hulls, by Proposition~\ref{prop:response}(a). The modular
operators are restrictions of the ambient operator to different vacuum-cyclic spaces.
The equality applies to excitations lying in every hull; arbitrary states can have
different relative entropies on those spaces.
Relative entropy can differ: the hull of an energy reader underestimates the information in a current
excitation that the hull of a current reader sees in full. Whatever role a selected
algebra plays in a gravitational response, its dependence on the seed enters through
excitations outside the Virasoro algebra; stress excitations alone cannot detect which
seed was selected. In higher dimensions Consciousness~II shows that selected
algebras with the required stress and boundary readings, on a conformal
net in Minkowski space and conditional on its supplied
$\mathrm{MC}_\eta$ ball-consistency law and source regularity, give the same
matter-sourced linearized Einstein tensor when they use one source and normalization; their relative entropies can differ~\cite{PaperCII}. That paper also constructs an exact neutral scalar preparation
inside a proper stable hull: a seed containing the stress readings and one invariant
quadratic observable contains a local unitary that produces nonzero Fourier weight.
With permission to implement this pulse, the observer can test the intrinsic response
without a charged-field phase reference. This construction is in the higher-dimensional
matter theory; the chiral Cut classification here does not supply its spacetime or
its modular-consistency law.

\section{What is still supplied, and three conjectures}\label{sec:open}

The construction derives the fixed algebra, its expectation and its modular invariance
from stationarity; the hull and its expectation from a seed; a proper Cut in every
nonscalar translation-stable hull from an ambient positive translation structure with a unique invariant vector;
and, under the hypotheses of Theorem~\ref{thm:galois}, the hull from the seed's
stabilizer. The following remain inputs.
\begin{enumerate}[label=(\arabic*),leftmargin=2em]
\item The ambient algebra, its state and, for Cuts, its translations. Hulls inherit
  Cuts; they do not create the ambient one. An inclusion of two hulls is never a Cut,
  since both are invariant under the modular group in both directions; one-sided
  objects such as the past modular orbit of Consciousness~I are needed for that.
\item The coupling to the environment. HAFF-B derives the selected algebra from the
  couplings, and Consciousness~I from the noise group; neither derives the couplings.
\item Stationarity, or a stationary partner supplied by an invariant
  expectation (Proposition~\ref{prop:asymptotic}). Relaxation requires
  an additional dynamical statement (Example~\ref{ex:mean}).
\item The task that fixes the foreground. T-DOME~II gives the optimal reading directions
  for a given task; the task is supplied.
\item A minimal compact symmetry for the classification. In the SU(2)\(_1\) net it is
  verified; elsewhere the hull is still determined by Theorems~\ref{thm:hull}
  and~\ref{thm:stableseed}, but its identification with a fixed-point algebra is open.
\end{enumerate}

\begin{conjecture}[Quantum symmetry fixes the hull]\label{conj:quantum}
Let \(\mathcal A\) be a completely rational conformal net and \(\mathcal B\subset\mathcal A\)
a M\"obius covariant subnet of finite index with \(\mathcal B(I)'\cap\mathcal A(I)=\mathbb C\mathbf 1\).
For every seed whose \(\Gamma_3\)-hull contains \(\mathcal B(0,\infty)\), the hull is the
half-line algebra of the intermediate net fixed by the generalized stabilizer of the
seed in the quantum symmetry of \(\mathcal B\subset\mathcal A\), and its vacuum character
is the corresponding partial sum of branching functions.
\end{conjecture}

The mechanism is Theorem~\ref{thm:galois} with the compact group replaced by the
symmetry that a finite-index inclusion carries. Support comes from the compact case
proved here, from the extension of the Izumi--Longo--Popa correspondence to compact Kac
algebras~\cite{ILP1998}, and from the correspondence between intermediate extensions
and intermediate Q-systems~\cite{BKLR2015}. The strongest objection is that a
\(\Gamma_3\)-hull is invariant under dilations and forward translations but is not
known in general to be the half-line algebra of a local subnet; in the SU(2)\(_1\) case
this follows from the Galois law, not from the hull construction. Without the trivial
relative commutant the conjecture fails, since intermediate algebras then need not be
fixed by any symmetry.

\begin{conjecture}[Window selection in type~III]\label{conj:window}
For the noise \(D^{H_z}_t\circ D^{SO(3)}_{st}\) and seeds containing a stress
observable, a window with \(s\kappa T\ll\varepsilon\ll\kappa T\) selects
\(M^{H_z}\) as the unique largest stable hull in this class: it is \(\varepsilon\)-stable on \([0,T]\), and no hull \(M^K\) with
\(K\not\supseteq H_z\) is. As the window grows past \(\varepsilon/(s\kappa)\) the selected
hull drops to \(\Vir_1(0,\infty)\), and the step is visible in the vacuum character.
\end{conjecture}

This is HAFF-B's window theorem~\cite{HAFFB} in type~III. Proposition~\ref{prop:plateau}
proves the stability half. For \(K\subsetneq H_z\), an element \(x\in M^K\) with
\(Q_{H_z}(x)=0\) and \(\|x\|=1\) has \(\|\mathcal E_t(x)-x\|=1-e^{-(1+s)\kappa t}\), which
proves the instability half for those \(K\). The objection is that for other \(K\) the
supremum of the instability over the unit ball of \(M^K\) has no evident lower bound
independent of \(K\); localized charged unitaries are the natural candidates for
proving one.

\begin{conjecture}[Every hull contains the Virasoro algebra]\label{conj:vir}
In the SU(2)\(_1\) net every nonscalar \(\Gamma_3\)-hull contains \(\Vir_1(0,\infty)\).
Then Theorem~\ref{thm:su2} classifies all hulls of the net.
\end{conjecture}

The mechanism is Theorem~\ref{thm:cut}: a nonscalar hull carries a Cut, and if the Cut
is standard it reconstructs a conformal subnet \(\mathcal B\). If \(\mathcal B\) is
diffeomorphism covariant with central charge one, the coset of \(\mathcal B\) in
\(\mathcal A\) has central charge zero and is trivial, so the stress tensor of
\(\mathcal B\) is the Sugawara tensor and \(\Vir_1\subseteq\mathcal B\)~\cite{GKO1986}.
A M\"obius covariant subnet of central charge below one cannot occur: Carpi showed that every M\"obius covariant subsystem of the SU(2)\(_1\) net is a fixed-point net \(\mathcal A^H\) for a closed subgroup \(H\subseteq SO(3)\), and all of these contain \(\Vir_1\)~\cite{Carpi1999}. The conjecture therefore reduces to a lifting problem, which is the objection: the net reconstructed from a standard hull Cut must be shown to be a conformal subsystem of the ambient net, that is, its local algebras \(\mathcal H\cap(U(1)\mathcal HU(1)^*)'\) must lie in the ambient local algebras and its M\"obius representation must be the restriction of the ambient one. Standardness of the inherited Cut is also not automatic.

\section{Conclusion}\label{sec:conclusion}

The physical reader and the ideal closure carry different information.
A nonzero-gain probe can generate a sharp pointer algebra without being
able to read that pointer sharply in one shot. In the explicit
nondemolition model, repetition with fresh ports makes the ideal
instrument available to any prescribed finite accuracy, with a budget
depending on gain. The instrument-level estimate survives old
references and finite adaptive composition. This is the task bridge;
spectral closure alone cannot supply it.

Stationarity of the source makes the fixed algebra a modular-invariant expected
subalgebra and supplies a persistent projection with no memory
kernel, without requiring relaxation of every trajectory. The stable hull
of a seed fixed pointwise by the noise carries the source's own modular data and lies
between the seed and the fixed algebra. A nonscalar hull that is also closed under
forward translations carries a proper Cut on its vacuum space when the ambient
positive translation group leaves only the vacuum invariant. Translation covariance
of the noise keeps this larger hull in the fixed algebra. Under a minimal compact
symmetry, a hull containing the symmetry-fixed algebra is fixed by the stabilizer of
the seed when its closure maps preserve the stabilizer-fixed algebra. Within this
class the accessible algebra is unique relative to what is read and changes only
when the symmetry left unbroken by the seed changes. In the SU(2)\(_1\) net the
vacuum characters of the hulls tell the stabilizers apart; a single current reading
generates the U(1) current algebra under the modular flow. For the
positive log-bump stress class, modular flow alone generates the Virasoro
half-line algebra, and the one-sided orbit constructs a proper standard
Cut on its own vacuum space. In the stated stress/current seed class,
dephasing about one axis leaves exactly three
ideal hulls, and each lasts as long as the noise respects
its symmetry. The response of a hull to stress excitations does not depend on the
seed, while its response to a current excitation does.

The source--access--Cut chain has a derived net-theoretic middle. The retained
seed fixes the ideal algebra through symmetry in the classified cases;
physical protocols decide which distinctions can be used within a task
budget. The translation-stable algebra
inherits the source's modular data and a Cut. In the SU(2)\(_1\) realization that Cut
is standard and reconstructs its conformal subnet on the inherited light ray.
Different seeds can select nonisomorphic nets there without producing distinct spacetime
metrics. The ambient net and physical reader remain inputs. The log-bump
construction removes an independent forward-translation requirement
from ideal stress-seed generation, but still presupposes the ambient
stress tensor and modular evolution. Coupling, preparation permissions
and the task remain physical inputs. Learned access and retained-history mechanisms
can change these inputs during an actual run; the hull theorem then
identifies the structure supported by the resulting seed. It does not
by itself explain the origin of the first available coupling.

\paragraph{Centred and uncentred task losses.}
A rank optimizer for centred omission loss need not optimize an
uncentred squared loss: the latter also contains
$(\sum_{i\notin I}\mathbb E\lambda_i)^2$.
For three independent channels with means and variances
$(0,1),(1,.1),(1,.1)$ and one retained channel, the uncentred minimizer
keeps a channel with smaller centred weight. The coding result in
Appendix~\ref{app:budget} therefore retains its stated task assumptions.

\appendix

\section{Rotation angles and conjugacy classes}\label{app:angles}

\begin{lemma}\label{lem:angles}
For a closed subgroup \(K\subseteq SO(3)\) let \(\mu_K\) be the distribution of the
rotation angle \(\theta\in[0,\pi]\) under normalized Haar measure. Then \(\mu_K\)
determines \(K\) up to conjugacy.
\end{lemma}

\begin{proof}
Up to conjugacy the closed subgroups are \(SO(3)\), \(O(2)\), \(SO(2)\), the
icosahedral, octahedral and tetrahedral groups \(I,O,T\) of orders \(60,24,12\), the
dihedral groups \(D_n\) of order \(2n\), \(n\ge2\), and the cyclic groups \(C_n\),
\(n\ge1\). The three infinite groups have \(\mu_K\) with density
\((1-\cos\theta)/\pi\), the density \(1/\pi\), and half the density \(1/\pi\) plus an
atom \(\frac12\) at \(\pi\); these correspond respectively to \(SO(3)\),
\(SO(2)\), and \(O(2)\). The finite ones are atomic. For finite \(K\) the atom at
\(0\) is \(1/|K|\), which fixes the order. Among groups of equal order, the atom at
\(\pi\) separates them: it is \(1/n\) for \(C_n\) with \(n\) even and \(0\) for \(n\)
odd; \((n+1)/2n\) for \(D_n\) with \(n\) even and \(\frac12\) for \(n\) odd; \(\frac14\)
for \(T\), \(\frac38\) for \(O\) and \(\frac14\) for \(I\). For order \(2m\), \(C_{2m}\)
has \(1/2m\) and \(D_m\) at least \(\frac12\), equal only for \(m=1\), where
\(D_1\) is conjugate to \(C_2\). For orders \(12\), \(24\) and \(60\) the three values
are \((\frac1{12},\frac7{12},\frac14)\), \((\frac1{24},\frac{13}{24},\frac38)\) and
\((\frac1{60},\frac{31}{60},\frac14)\) for \(C,D,\) and \(T,O,I\) respectively.
\end{proof}

\section{Numerical checks}\label{app:checks}

The supplementary source map identifies the following diagnostics
and their retained outputs. Finite samples test the stated identities;
the infinite-dimensional conclusions use the proofs in the text.

\paragraph{B.0 Physical permission bridge.}
The controlled-X complete instrument is checked on an old Bell reference
and eleven random mixed inputs at $\alpha=0,.1,.6,1$ and $m=1,2,3,5$,
including the fair tie rule and post-read-flip counterexample.
The sharp half-diamond identity is analytic
(Proposition~\ref{prop:readbridge}).

\paragraph{B.1 Finite hulls.}
Closure under products, adjoints, the modular generator and noise
reproduces the nonstationary qubit hull $M_2$, its pinched-state
two-dimensional hull, and the $M_2\otimes1$ stationary fixed sector.
The three-qubit example has seed, hull and fixed dimensions
$3,10,20$, respectively. Other checks distinguish equal type-I fixed
algebras with different stabilizers, verify persistence and joins,
and test Proposition~\ref{prop:alignment} on 200 task matrices.

\paragraph{B.2 Characters.}
The class data of $C_n,D_n$ for $n\le40$, the three finite polyhedral
groups and $SO(2),O(2),SO(3)$ give 85 distinct multiplicity sequences
through $j=60$. Their characters reproduce the table and the
Virasoro, current, SU(2) level-one and even-Fock low-energy coefficients.
Finite-level agreement does not replace the analytic character
and rotation-angle arguments.

\section{A separate finite-coordinate coding model}\label{app:budget}

Here $V$ is an abstract finite-dimensional task coordinate space.
The propositions do not construct a net observable from a coordinate
projection and are not premises of Corollary~\ref{cor:onechannel}.
T-DOME~II supplies the Gaussian task optimizer used in this model.

\begin{proposition}[Rank under a processing budget]\label{prop:budget}
Let a reader hold a real coordinate space \(V\) of dimension \(D\), let admissible
encoders be orthogonal projections \(Q\) on \(V\) with processing cost
\(h_\mu\operatorname{rank}Q\), \(h_\mu>0\), and assume:
\begin{enumerate}[label=(\arabic*),leftmargin=2em]
\item the full-rank encoder is infeasible, \(h_\mu D>\dot{\mathcal C}_{\mathrm{budget}}\);
\item the distortion functional attains its minimum on the feasible projections and
  strictly decreases whenever a feasible projection is enlarged by one coordinate;
\item a group \(G\) acts irreducibly and orthogonally on \(V\), the distortion is
  \(G\)-invariant, and \(0<\lfloor\dot{\mathcal C}_{\mathrm{budget}}/h_\mu\rfloor<D\).
\end{enumerate}
Then every minimizing encoder has rank
\(k^*=\lfloor\dot{\mathcal C}_{\mathrm{budget}}/h_\mu\rfloor<D\), its range is a
\(k^*\)-dimensional foreground, and no such projection is \(G\)-invariant, so choosing a
minimizer breaks the \(G\)-symmetry.
\end{proposition}

\begin{proof}
Feasibility gives \(\operatorname{rank}Q\le k^*\), and strict improvement excludes
smaller rank. A minimizing projection that commuted with the irreducible representation
would have invariant range and kernel, hence would be \(0\) or \(\mathbf 1\), contrary
to \(0<k^*<D\).
\end{proof}

The rank charge, strict improvement and irreducible action are inputs of this coding
model. They are not consequences of persistence, and the proposition does not say that
symmetry breaking is unavoidable outside this class of encoders.

Which \(k^*\) directions are kept is answered by T-DOME~II. For a budgeted reader whose
task is linear in a fresh Gaussian input and whose loss is quadratic, with
all rank-\(k\) linear projectors allowed, the best risk of a rank-\(k\) reading policy is
attained by the top-\(k\) eigenspace of the posterior task matrix \(\widehat W_M\)
built from its memory \(M\), and the optimal subspace is unique when
\(\lambda_k(\widehat W_M)>\lambda_{k+1}(\widehat W_M)\)~\cite{TDOMEII}. These directions are
record channels when \(\widehat W_M\) commutes with the projections onto the record
channels, and approximately so otherwise.

\begin{proposition}[Foreground and record channels]\label{prop:alignment}
Let the reader's \(D\) channels be record channels, and write
\(\widehat W_M=W_{\mathrm{d}}+W_{\mathrm{o}}\) with \(W_{\mathrm d}\) diagonal in the
channel basis and \(W_{\mathrm o}\) its off-diagonal part. Let \(\Delta_k>0\) be the gap
between the \(k\)-th and \((k+1)\)-st largest diagonal entries and \(P_{\mathrm{ch}}\)
the projection onto the corresponding \(k\) channels. If \(W_{\mathrm o}=0\), the
optimal foreground is spanned by those \(k\) channels. If
\(\|W_{\mathrm o}\|<\Delta_k\), any rank-\(k\) optimal projection \(P_{\mathrm{opt}}\) satisfies
\[
\|P_{\mathrm{opt}}-P_{\mathrm{ch}}\|\le\frac{\|W_{\mathrm o}\|}{\Delta_k-\|W_{\mathrm o}\|}.
\]
\end{proposition}

\begin{proof}
The first statement is the capability theorem of T-DOME~II for a diagonal matrix. For the
second, Weyl's inequality separates the top \(k\) eigenvalues of \(\widehat W_M\) from
the remaining diagonal entries of \(W_{\mathrm d}\) by at least
\(\Delta_k-\|W_{\mathrm o}\|\), and the Davis--Kahan \(\sin\Theta\)
theorem~\cite{DavisKahan1970} bounds the sine of the largest principal angle, which for
projections of equal rank equals \(\|P_{\mathrm{opt}}-P_{\mathrm{ch}}\|\), by
\(\|W_{\mathrm o}\|\) divided by that separation.
\end{proof}

\clearpage
\addcontentsline{toc}{section}{References}

\begin{thebibliography}{99}

\bibitem{PaperCII}
S.~Liu, \emph{Higher-Dimensional Consistency of Geometry and Modular Flow: Physical Source Identification and a Second-Order Extension} (Consciousness~II), (2026), series record DOI: 10.5281/zenodo.19607543.
\bibitem{PaperI}
S.~Liu, \emph{Geometry Relative to a Cut: Selection, Generation and Modular
Inseparability}, 2026, July series record DOI:
10.5281/zenodo.19607387.

\bibitem{HAFFA}
S.~Liu, \emph{Physical Access and Correlation Geometry: The Holographic
Alaya--Field Framework}, (2026), series record DOI:
10.5281/zenodo.18361706.

\bibitem{HAFFB}
S.~Liu, \emph{Coupling, Records, and Stable Relations:
Formation, Protected Content and the Lifetime of a Geometry}, (2026), series record
DOI: 10.5281/zenodo.18367060.

\bibitem{HAFFD}
S.~Liu, \emph{Selection, Maintenance and Gravitational Response: Modular Constraints and a Physical Return}, (2026), series
record DOI: 10.5281/zenodo.18388881.

\bibitem{HAFFE}
S.~Liu, \emph{Measurement as Accessibility: Measurement Permissions, Stable Records and
Public Readout}, (2026), series record DOI: 10.5281/zenodo.18400065.

\bibitem{HAFFF}
S.~Liu, \emph{Records, Dissipation and Intervention: Conditions for a Temporal
Direction}, (2026), series record DOI: 10.5281/zenodo.18400425.

\bibitem{TDOMEII}
S.~Liu, \emph{Reference Frames Selected by History: The Capability and Cost of a Bounded
Agent's Self-Frame}, (2026), series record DOI: 10.5281/zenodo.18579703.

\bibitem{QRAIFB}
S.~Liu, \emph{Thermodynamic Stability Constraints on Persistent Open Quantum Subsystems:
Control Algebra, Physical Actuation and Energy Cost}, (2026), series
record DOI: 10.5281/zenodo.18525890.

\bibitem{AZ}
H.~Araki and L.~Zsid\'o, Extension of the structure theorem of Borchers and its
application to half-sided modular inclusions, \emph{Rev.\ Math.\ Phys.} \textbf{17}
(2005) 491--543.

\bibitem{BKLR2015}
M.~Bischoff, Y.~Kawahigashi, R.~Longo and K.-H.~Rehren, \emph{Tensor Categories and
Endomorphisms of von Neumann Algebras, with Applications to Quantum Field Theory},
SpringerBriefs in Mathematical Physics~3, Springer (2015).

\bibitem{Borchers1992}
H.-J.~Borchers, The CPT-theorem in two-dimensional theories of local observables,
\emph{Commun.\ Math.\ Phys.} \textbf{143} (1992) 315--332.

\bibitem{BGL}
R.~Brunetti, D.~Guido and R.~Longo, Modular structure and duality in conformal quantum
field theory, \emph{Commun.\ Math.\ Phys.} \textbf{156} (1993) 201--219.

\bibitem{Carpi1999}
S.~Carpi, Classification of subsystems for the Haag--Kastler nets generated by \(c=1\) chiral current algebras, \emph{Lett.\ Math.\ Phys.} \textbf{47} (1999) 353--364.
\bibitem{Carpi}
S.~Carpi, On the representation theory of Virasoro nets, \emph{Commun.\ Math.\ Phys.}
\textbf{244} (2004) 261--284.

\bibitem{Choi1974}
M.-D.~Choi, A Schwarz inequality for positive linear maps on C*-algebras,
\emph{Illinois J.\ Math.} \textbf{18} (1974) 565--574.

\bibitem{DavisKahan1970}
C.~Davis and W.~M.~Kahan, The rotation of eigenvectors by a perturbation.~III,
\emph{SIAM J.\ Numer.\ Anal.} \textbf{7} (1970) 1--46.

\bibitem{GKO1986}
P.~Goddard, A.~Kent and D.~Olive, Unitary representations of the Virasoro and
super-Virasoro algebras, \emph{Commun.\ Math.\ Phys.} \textbf{103} (1986) 105--119.

\bibitem{GLW}
D.~Guido, R.~Longo and H.-W.~Wiesbrock, Extensions of conformal nets and superselection
structures, \emph{Commun.\ Math.\ Phys.} \textbf{192} (1998) 217--244.

\bibitem{Hollands}
S.~Hollands, Relative entropy for coherent states in chiral CFT, arXiv:1903.07508.

\bibitem{ILP1998}
M.~Izumi, R.~Longo and S.~Popa, A Galois correspondence for compact groups of
automorphisms of von Neumann algebras with a generalization to Kac algebras,
\emph{J.\ Funct.\ Anal.} \textbf{155} (1998) 25--63.

\bibitem{KR}
V.~G.~Kac and A.~K.~Raina, \emph{Highest Weight Representations of Infinite Dimensional
Lie Algebras}, World Scientific (1987).

\bibitem{Watanabe1979}
S.~Watanabe, Ergodic theorems for dynamical semi-groups on operator algebras, \emph{Hokkaido Math.\ J.} \textbf{8} (1979) 176--190.
\bibitem{KovacsSzucs1966}
I.~Kov\'acs and J.~Sz\H{u}cs, Ergodic type theorems in von Neumann algebras, \emph{Acta
Sci.\ Math.\ (Szeged)} \textbf{27} (1966) 233--246.

\bibitem{JencovaPetz2006}
A.~Jen\v{c}ov\'a and D.~Petz, Sufficiency in quantum statistical inference: a survey with examples, arXiv:quant-ph/0604091 (2006).
\bibitem{Petz1986}
D.~Petz, Sufficient subalgebras and the relative entropy of states of a von Neumann
algebra, \emph{Commun.\ Math.\ Phys.} \textbf{105} (1986) 123--131.

\bibitem{Rehren}
K.-H.~Rehren, A new view of the Virasoro algebra, \emph{Lett.\ Math.\ Phys.} \textbf{30}
(1994) 125--130.

\bibitem{Segal}
G.~Segal, Unitary representations of some infinite dimensional groups, \emph{Commun.\
Math.\ Phys.} \textbf{80} (1981) 301--342.

\bibitem{Takesaki1972}
M.~Takesaki, Conditional expectations in von Neumann algebras, \emph{J.\ Funct.\ Anal.}
\textbf{9} (1972) 306--321.

\bibitem{Wiener1932}
N.~Wiener, Tauberian theorems, \emph{Ann.\ of Math.} \textbf{33} (1932) 1--100.

\bibitem{Wiesbrock1993}
H.-W.~Wiesbrock, Half-sided modular inclusions of von Neumann algebras, \emph{Commun.\
Math.\ Phys.} \textbf{157} (1993) 83--92.

\bibitem{Xu}
F.~Xu, Strong additivity and conformal nets, \emph{Pacific J.\ Math.} \textbf{221}
(2005) 167--199.

\end{thebibliography}

\end{document}
